% Options for packages loaded elsewhere
\PassOptionsToPackage{unicode}{hyperref}
\PassOptionsToPackage{hyphens}{url}
\documentclass[
]{article}
\usepackage{xcolor}
\usepackage[margin=1in]{geometry}
\usepackage{amsmath,amssymb}
\setcounter{secnumdepth}{-\maxdimen} % remove section numbering
\usepackage{iftex}
\ifPDFTeX
  \usepackage[T1]{fontenc}
  \usepackage[utf8]{inputenc}
  \usepackage{textcomp} % provide euro and other symbols
\else % if luatex or xetex
  \usepackage{unicode-math} % this also loads fontspec
  \defaultfontfeatures{Scale=MatchLowercase}
  \defaultfontfeatures[\rmfamily]{Ligatures=TeX,Scale=1}
\fi
\usepackage{lmodern}
\ifPDFTeX\else
  % xetex/luatex font selection
\fi
% Use upquote if available, for straight quotes in verbatim environments
\IfFileExists{upquote.sty}{\usepackage{upquote}}{}
\IfFileExists{microtype.sty}{% use microtype if available
  \usepackage[]{microtype}
  \UseMicrotypeSet[protrusion]{basicmath} % disable protrusion for tt fonts
}{}
\makeatletter
\@ifundefined{KOMAClassName}{% if non-KOMA class
  \IfFileExists{parskip.sty}{%
    \usepackage{parskip}
  }{% else
    \setlength{\parindent}{0pt}
    \setlength{\parskip}{6pt plus 2pt minus 1pt}}
}{% if KOMA class
  \KOMAoptions{parskip=half}}
\makeatother
\usepackage{color}
\usepackage{fancyvrb}
\newcommand{\VerbBar}{|}
\newcommand{\VERB}{\Verb[commandchars=\\\{\}]}
\DefineVerbatimEnvironment{Highlighting}{Verbatim}{commandchars=\\\{\}}
% Add ',fontsize=\small' for more characters per line
\usepackage{framed}
\definecolor{shadecolor}{RGB}{248,248,248}
\newenvironment{Shaded}{\begin{snugshade}}{\end{snugshade}}
\newcommand{\AlertTok}[1]{\textcolor[rgb]{0.94,0.16,0.16}{#1}}
\newcommand{\AnnotationTok}[1]{\textcolor[rgb]{0.56,0.35,0.01}{\textbf{\textit{#1}}}}
\newcommand{\AttributeTok}[1]{\textcolor[rgb]{0.13,0.29,0.53}{#1}}
\newcommand{\BaseNTok}[1]{\textcolor[rgb]{0.00,0.00,0.81}{#1}}
\newcommand{\BuiltInTok}[1]{#1}
\newcommand{\CharTok}[1]{\textcolor[rgb]{0.31,0.60,0.02}{#1}}
\newcommand{\CommentTok}[1]{\textcolor[rgb]{0.56,0.35,0.01}{\textit{#1}}}
\newcommand{\CommentVarTok}[1]{\textcolor[rgb]{0.56,0.35,0.01}{\textbf{\textit{#1}}}}
\newcommand{\ConstantTok}[1]{\textcolor[rgb]{0.56,0.35,0.01}{#1}}
\newcommand{\ControlFlowTok}[1]{\textcolor[rgb]{0.13,0.29,0.53}{\textbf{#1}}}
\newcommand{\DataTypeTok}[1]{\textcolor[rgb]{0.13,0.29,0.53}{#1}}
\newcommand{\DecValTok}[1]{\textcolor[rgb]{0.00,0.00,0.81}{#1}}
\newcommand{\DocumentationTok}[1]{\textcolor[rgb]{0.56,0.35,0.01}{\textbf{\textit{#1}}}}
\newcommand{\ErrorTok}[1]{\textcolor[rgb]{0.64,0.00,0.00}{\textbf{#1}}}
\newcommand{\ExtensionTok}[1]{#1}
\newcommand{\FloatTok}[1]{\textcolor[rgb]{0.00,0.00,0.81}{#1}}
\newcommand{\FunctionTok}[1]{\textcolor[rgb]{0.13,0.29,0.53}{\textbf{#1}}}
\newcommand{\ImportTok}[1]{#1}
\newcommand{\InformationTok}[1]{\textcolor[rgb]{0.56,0.35,0.01}{\textbf{\textit{#1}}}}
\newcommand{\KeywordTok}[1]{\textcolor[rgb]{0.13,0.29,0.53}{\textbf{#1}}}
\newcommand{\NormalTok}[1]{#1}
\newcommand{\OperatorTok}[1]{\textcolor[rgb]{0.81,0.36,0.00}{\textbf{#1}}}
\newcommand{\OtherTok}[1]{\textcolor[rgb]{0.56,0.35,0.01}{#1}}
\newcommand{\PreprocessorTok}[1]{\textcolor[rgb]{0.56,0.35,0.01}{\textit{#1}}}
\newcommand{\RegionMarkerTok}[1]{#1}
\newcommand{\SpecialCharTok}[1]{\textcolor[rgb]{0.81,0.36,0.00}{\textbf{#1}}}
\newcommand{\SpecialStringTok}[1]{\textcolor[rgb]{0.31,0.60,0.02}{#1}}
\newcommand{\StringTok}[1]{\textcolor[rgb]{0.31,0.60,0.02}{#1}}
\newcommand{\VariableTok}[1]{\textcolor[rgb]{0.00,0.00,0.00}{#1}}
\newcommand{\VerbatimStringTok}[1]{\textcolor[rgb]{0.31,0.60,0.02}{#1}}
\newcommand{\WarningTok}[1]{\textcolor[rgb]{0.56,0.35,0.01}{\textbf{\textit{#1}}}}
\usepackage{longtable,booktabs,array}
\usepackage{calc} % for calculating minipage widths
% Correct order of tables after \paragraph or \subparagraph
\usepackage{etoolbox}
\makeatletter
\patchcmd\longtable{\par}{\if@noskipsec\mbox{}\fi\par}{}{}
\makeatother
% Allow footnotes in longtable head/foot
\IfFileExists{footnotehyper.sty}{\usepackage{footnotehyper}}{\usepackage{footnote}}
\makesavenoteenv{longtable}
\usepackage{graphicx}
\makeatletter
\newsavebox\pandoc@box
\newcommand*\pandocbounded[1]{% scales image to fit in text height/width
  \sbox\pandoc@box{#1}%
  \Gscale@div\@tempa{\textheight}{\dimexpr\ht\pandoc@box+\dp\pandoc@box\relax}%
  \Gscale@div\@tempb{\linewidth}{\wd\pandoc@box}%
  \ifdim\@tempb\p@<\@tempa\p@\let\@tempa\@tempb\fi% select the smaller of both
  \ifdim\@tempa\p@<\p@\scalebox{\@tempa}{\usebox\pandoc@box}%
  \else\usebox{\pandoc@box}%
  \fi%
}
% Set default figure placement to htbp
\def\fps@figure{htbp}
\makeatother
\setlength{\emergencystretch}{3em} % prevent overfull lines
\providecommand{\tightlist}{%
  \setlength{\itemsep}{0pt}\setlength{\parskip}{0pt}}
\usepackage{bookmark}
\IfFileExists{xurl.sty}{\usepackage{xurl}}{} % add URL line breaks if available
\urlstyle{same}
\hypersetup{
  pdftitle={Agrupamiento de Titulaciones en Clústeres: Análisis de Rendimiento y Abandono},
  pdfauthor={Equipo Data Science},
  hidelinks,
  pdfcreator={LaTeX via pandoc}}

\title{Agrupamiento de Titulaciones en Clústeres: Análisis de
Rendimiento y Abandono}
\author{Equipo Data Science}
\date{2026-03-03}

\begin{document}
\maketitle

{
\setcounter{tocdepth}{2}
\tableofcontents
}
\section{1. Introducción}\label{introducciuxf3n}

Este análisis responde a la pregunta grupal: \textbf{``Las titulaciones
pueden agruparse en clústeres definidos que faciliten la intervención
institucional''}

\subsection{1.1 Hipótesis}\label{hipuxf3tesis}

\begin{itemize}
\tightlist
\item
  \textbf{H₀}: No es posible identificar grupos de titulaciones con
  comportamientos homogéneos de bajo rendimiento y alto abandono.
\item
  \textbf{H₁}: Las titulaciones pueden agruparse en clústeres definidos
  que faciliten la intervención institucional (titulaciones en riesgo
  vs.~exitosas).
\end{itemize}

\subsection{1.2 Diseño del Análisis}\label{diseuxf1o-del-anuxe1lisis}

\begin{itemize}
\tightlist
\item
  \textbf{Variables}: Tasa de Rendimiento (1.6) y Tasa de Abandono
  (1.11)
\item
  \textbf{Método}: K-means clustering tras normalización de datos
\item
  \textbf{Visualización}: Gráfico de dispersión coloreado por clúster
\end{itemize}

\begin{center}\rule{0.5\linewidth}{0.5pt}\end{center}

\section{2. Carga y Exploración de
Datos}\label{carga-y-exploraciuxf3n-de-datos}

\begin{Shaded}
\begin{Highlighting}[]
\CommentTok{\# Cargar datos}
\NormalTok{df }\OtherTok{\textless{}{-}} \FunctionTok{read\_csv}\NormalTok{(}\StringTok{"data/dataset\_abandono\_rendimiento\_estricto.csv"}\NormalTok{)}

\CommentTok{\# Visualizar estructura de los datos}
\FunctionTok{str}\NormalTok{(df)}
\end{Highlighting}
\end{Shaded}

\begin{verbatim}
## spc_tbl_ [756 x 10] (S3: spec_tbl_df/tbl_df/tbl/data.frame)
##  $ Cod_Centro              : num [1:756] 1 1 1 1 1 1 1 5 5 5 ...
##  $ Cod_Titulacion          : chr [1:756] "151" "151" "151" "151" ...
##  $ Titulacion              : chr [1:756] "Grado en Administración y Dirección de Empresas" "Grado en Administración y Dirección de Empresas" "Grado en Administración y Dirección de Empresas" "Grado en Administración y Dirección de Empresas" ...
##  $ Año                     : chr [1:756] "2017_2018" "2018_2019" "2019_2020" "2020_2021" ...
##  $ Tasa_Rendimiento        : num [1:756] 0.553 0.57 0.685 0.579 0.533 ...
##  $ Alumnos_Egresados       : num [1:756] 63348 64920 74334 65244 59856 ...
##  $ Matriculados_Rendimiento: num [1:756] 114624 113796 108534 112680 112326 ...
##  $ Tasa_Abandono           : num [1:756] 0.311 0.365 0.371 0.387 0.375 ...
##  $ Alumnos_Abandonan       : num [1:756] 160 183 182 191 187 200 182 33 37 37 ...
##  $ Matriculados_Abandono   : num [1:756] 514 501 491 493 499 492 480 220 224 212 ...
##  - attr(*, "spec")=
##   .. cols(
##   ..   Cod_Centro = col_double(),
##   ..   Cod_Titulacion = col_character(),
##   ..   Titulacion = col_character(),
##   ..   Año = col_character(),
##   ..   Tasa_Rendimiento = col_double(),
##   ..   Alumnos_Egresados = col_double(),
##   ..   Matriculados_Rendimiento = col_double(),
##   ..   Tasa_Abandono = col_double(),
##   ..   Alumnos_Abandonan = col_double(),
##   ..   Matriculados_Abandono = col_double()
##   .. )
##  - attr(*, "problems")=<externalptr>
\end{verbatim}

\begin{Shaded}
\begin{Highlighting}[]
\CommentTok{\# Primeras filas}
\FunctionTok{head}\NormalTok{(df, }\DecValTok{10}\NormalTok{)}
\end{Highlighting}
\end{Shaded}

\begin{verbatim}
## # A tibble: 10 x 10
##    Cod_Centro Cod_Titulacion Titulacion Año   Tasa_Rendimiento Alumnos_Egresados
##         <dbl> <chr>          <chr>      <chr>            <dbl>             <dbl>
##  1          1 151            Grado en ~ 2017~            0.553             63348
##  2          1 151            Grado en ~ 2018~            0.570             64920
##  3          1 151            Grado en ~ 2019~            0.685             74334
##  4          1 151            Grado en ~ 2020~            0.579             65244
##  5          1 151            Grado en ~ 2021~            0.533             59856
##  6          1 151            Grado en ~ 2022~            0.547             63084
##  7          1 151            Grado en ~ 2023~            0.586             69042
##  8          5 153            Grado en ~ 2017~            0.768             45096
##  9          5 153            Grado en ~ 2018~            0.751             41418
## 10          5 153            Grado en ~ 2019~            0.803             43152
## # i 4 more variables: Matriculados_Rendimiento <dbl>, Tasa_Abandono <dbl>,
## #   Alumnos_Abandonan <dbl>, Matriculados_Abandono <dbl>
\end{verbatim}

\begin{Shaded}
\begin{Highlighting}[]
\CommentTok{\# Resumen de estadísticas descriptivas}
\FunctionTok{summary}\NormalTok{(df)}
\end{Highlighting}
\end{Shaded}

\begin{verbatim}
##    Cod_Centro     Cod_Titulacion      Titulacion            Año           
##  Min.   :  1.00   Length:756         Length:756         Length:756        
##  1st Qu.:  4.75   Class :character   Class :character   Class :character  
##  Median : 15.50   Mode  :character   Mode  :character   Mode  :character  
##  Mean   : 27.55                                                           
##  3rd Qu.: 21.00                                                           
##  Max.   :302.00                                                           
##  Tasa_Rendimiento Alumnos_Egresados Matriculados_Rendimiento Tasa_Abandono    
##  Min.   :0.1579   Min.   :    18    Min.   :   114           Min.   :0.00000  
##  1st Qu.:0.6304   1st Qu.:  1436    1st Qu.:  1816           1st Qu.:0.08522  
##  Median :0.7443   Median :  6342    Median : 11032           Median :0.15831  
##  Mean   :0.7315   Mean   : 14184    Mean   : 19904           Mean   :0.19837  
##  3rd Qu.:0.8499   3rd Qu.: 20687    3rd Qu.: 31436           3rd Qu.:0.29013  
##  Max.   :1.2016   Max.   :122221    Max.   :143226           Max.   :0.96203  
##  Alumnos_Abandonan Matriculados_Abandono
##  Min.   :  0.00    Min.   :  5.00       
##  1st Qu.:  3.00    1st Qu.: 27.75       
##  Median :  9.00    Median : 57.50       
##  Mean   : 22.02    Mean   : 98.12       
##  3rd Qu.: 32.00    3rd Qu.:137.00       
##  Max.   :211.00    Max.   :635.00
\end{verbatim}

\section{2.1 Información General del
Dataset}\label{informaciuxf3n-general-del-dataset}

\begin{Shaded}
\begin{Highlighting}[]
\FunctionTok{cat}\NormalTok{(}\StringTok{"Número de titulaciones únicas:"}\NormalTok{, }\FunctionTok{n\_distinct}\NormalTok{(df}\SpecialCharTok{$}\NormalTok{Titulacion), }\StringTok{"}\SpecialCharTok{\textbackslash{}n}\StringTok{"}\NormalTok{)}
\end{Highlighting}
\end{Shaded}

\begin{verbatim}
## Número de titulaciones únicas: 198
\end{verbatim}

\begin{Shaded}
\begin{Highlighting}[]
\FunctionTok{cat}\NormalTok{(}\StringTok{"Años disponibles:"}\NormalTok{, }\FunctionTok{paste}\NormalTok{(}\FunctionTok{unique}\NormalTok{(df}\SpecialCharTok{$}\NormalTok{Año), }\AttributeTok{collapse =} \StringTok{", "}\NormalTok{), }\StringTok{"}\SpecialCharTok{\textbackslash{}n}\StringTok{"}\NormalTok{)}
\end{Highlighting}
\end{Shaded}

\begin{verbatim}
## Años disponibles: 2017_2018, 2018_2019, 2019_2020, 2020_2021, 2021_2022, 2022_2023, 2023_2024
\end{verbatim}

\begin{Shaded}
\begin{Highlighting}[]
\FunctionTok{cat}\NormalTok{(}\StringTok{"Número total de observaciones:"}\NormalTok{, }\FunctionTok{nrow}\NormalTok{(df), }\StringTok{"}\SpecialCharTok{\textbackslash{}n}\StringTok{"}\NormalTok{)}
\end{Highlighting}
\end{Shaded}

\begin{verbatim}
## Número total de observaciones: 756
\end{verbatim}

\begin{Shaded}
\begin{Highlighting}[]
\FunctionTok{cat}\NormalTok{(}\StringTok{"Número de centros:"}\NormalTok{, }\FunctionTok{n\_distinct}\NormalTok{(df}\SpecialCharTok{$}\NormalTok{Cod\_Centro), }\StringTok{"}\SpecialCharTok{\textbackslash{}n}\StringTok{"}\NormalTok{)}
\end{Highlighting}
\end{Shaded}

\begin{verbatim}
## Número de centros: 24
\end{verbatim}

\begin{center}\rule{0.5\linewidth}{0.5pt}\end{center}

\section{3. Preparación de Datos}\label{preparaciuxf3n-de-datos}

\begin{Shaded}
\begin{Highlighting}[]
\CommentTok{\# Agregar datos por titulación (media de todos los años)}
\NormalTok{titulaciones\_agg }\OtherTok{\textless{}{-}}\NormalTok{ df }\SpecialCharTok{\%\textgreater{}\%}
  \FunctionTok{group\_by}\NormalTok{(Titulacion) }\SpecialCharTok{\%\textgreater{}\%}
  \FunctionTok{summarise}\NormalTok{(}
    \AttributeTok{Tasa\_Rendimiento =} \FunctionTok{mean}\NormalTok{(Tasa\_Rendimiento, }\AttributeTok{na.rm =} \ConstantTok{TRUE}\NormalTok{),}
    \AttributeTok{Tasa\_Abandono =} \FunctionTok{mean}\NormalTok{(Tasa\_Abandono, }\AttributeTok{na.rm =} \ConstantTok{TRUE}\NormalTok{),}
    \AttributeTok{n\_observaciones =} \FunctionTok{n}\NormalTok{(),}
    \AttributeTok{.groups =} \StringTok{\textquotesingle{}drop\textquotesingle{}}
\NormalTok{  ) }\SpecialCharTok{\%\textgreater{}\%}
  \CommentTok{\# Filtrar titulaciones con datos válidos}
  \FunctionTok{filter}\NormalTok{(}\SpecialCharTok{!}\FunctionTok{is.na}\NormalTok{(Tasa\_Rendimiento) }\SpecialCharTok{\&} \SpecialCharTok{!}\FunctionTok{is.na}\NormalTok{(Tasa\_Abandono))}

\FunctionTok{dim}\NormalTok{(titulaciones\_agg)}
\end{Highlighting}
\end{Shaded}

\begin{verbatim}
## [1] 198   4
\end{verbatim}

\begin{Shaded}
\begin{Highlighting}[]
\FunctionTok{head}\NormalTok{(titulaciones\_agg)}
\end{Highlighting}
\end{Shaded}

\begin{verbatim}
## # A tibble: 6 x 4
##   Titulacion                      Tasa_Rendimiento Tasa_Abandono n_observaciones
##   <chr>                                      <dbl>         <dbl>           <int>
## 1 Grado en Administración y Dire~            0.579         0.371               7
## 2 Grado en Antropología Social y~            0.625         0.452               7
## 3 Grado en Bellas Artes                      0.869         0.182               7
## 4 Grado en Biología                          0.761         0.157               7
## 5 Grado en Biomedicina Básica y ~            0.918         0.171               7
## 6 Grado en Bioquímica por la U. ~            0.869         0.136               2
\end{verbatim}

\begin{Shaded}
\begin{Highlighting}[]
\CommentTok{\# Visualizar la distribución de las variables}
\FunctionTok{summary}\NormalTok{(titulaciones\_agg[, }\FunctionTok{c}\NormalTok{(}\StringTok{"Tasa\_Rendimiento"}\NormalTok{, }\StringTok{"Tasa\_Abandono"}\NormalTok{)])}
\end{Highlighting}
\end{Shaded}

\begin{verbatim}
##  Tasa_Rendimiento Tasa_Abandono    
##  Min.   :0.3989   Min.   :0.00000  
##  1st Qu.:0.6598   1st Qu.:0.08263  
##  Median :0.7648   Median :0.12852  
##  Mean   :0.7458   Mean   :0.16771  
##  3rd Qu.:0.8439   3rd Qu.:0.23709  
##  Max.   :1.1334   Max.   :0.59175
\end{verbatim}

\section{3.1 Estadísticas Descriptivas por
Variable}\label{estaduxedsticas-descriptivas-por-variable}

\begin{Shaded}
\begin{Highlighting}[]
\FunctionTok{cat}\NormalTok{(}\StringTok{"TASA DE RENDIMIENTO}\SpecialCharTok{\textbackslash{}n}\StringTok{"}\NormalTok{)}
\end{Highlighting}
\end{Shaded}

\begin{verbatim}
## TASA DE RENDIMIENTO
\end{verbatim}

\begin{Shaded}
\begin{Highlighting}[]
\FunctionTok{cat}\NormalTok{(}\StringTok{"Media:"}\NormalTok{, }\FunctionTok{mean}\NormalTok{(titulaciones\_agg}\SpecialCharTok{$}\NormalTok{Tasa\_Rendimiento), }\StringTok{"}\SpecialCharTok{\textbackslash{}n}\StringTok{"}\NormalTok{)}
\end{Highlighting}
\end{Shaded}

\begin{verbatim}
## Media: 0.7458289
\end{verbatim}

\begin{Shaded}
\begin{Highlighting}[]
\FunctionTok{cat}\NormalTok{(}\StringTok{"Desv. Est.:"}\NormalTok{, }\FunctionTok{sd}\NormalTok{(titulaciones\_agg}\SpecialCharTok{$}\NormalTok{Tasa\_Rendimiento), }\StringTok{"}\SpecialCharTok{\textbackslash{}n}\StringTok{"}\NormalTok{)}
\end{Highlighting}
\end{Shaded}

\begin{verbatim}
## Desv. Est.: 0.1316349
\end{verbatim}

\begin{Shaded}
\begin{Highlighting}[]
\FunctionTok{cat}\NormalTok{(}\StringTok{"Mín:"}\NormalTok{, }\FunctionTok{min}\NormalTok{(titulaciones\_agg}\SpecialCharTok{$}\NormalTok{Tasa\_Rendimiento), }\StringTok{"}\SpecialCharTok{\textbackslash{}n}\StringTok{"}\NormalTok{)}
\end{Highlighting}
\end{Shaded}

\begin{verbatim}
## Mín: 0.3988735
\end{verbatim}

\begin{Shaded}
\begin{Highlighting}[]
\FunctionTok{cat}\NormalTok{(}\StringTok{"Máx:"}\NormalTok{, }\FunctionTok{max}\NormalTok{(titulaciones\_agg}\SpecialCharTok{$}\NormalTok{Tasa\_Rendimiento), }\StringTok{"}\SpecialCharTok{\textbackslash{}n\textbackslash{}n}\StringTok{"}\NormalTok{)}
\end{Highlighting}
\end{Shaded}

\begin{verbatim}
## Máx: 1.133439
\end{verbatim}

\begin{Shaded}
\begin{Highlighting}[]
\FunctionTok{cat}\NormalTok{(}\StringTok{"TASA DE ABANDONO}\SpecialCharTok{\textbackslash{}n}\StringTok{"}\NormalTok{)}
\end{Highlighting}
\end{Shaded}

\begin{verbatim}
## TASA DE ABANDONO
\end{verbatim}

\begin{Shaded}
\begin{Highlighting}[]
\FunctionTok{cat}\NormalTok{(}\StringTok{"Media:"}\NormalTok{, }\FunctionTok{mean}\NormalTok{(titulaciones\_agg}\SpecialCharTok{$}\NormalTok{Tasa\_Abandono), }\StringTok{"}\SpecialCharTok{\textbackslash{}n}\StringTok{"}\NormalTok{)}
\end{Highlighting}
\end{Shaded}

\begin{verbatim}
## Media: 0.1677057
\end{verbatim}

\begin{Shaded}
\begin{Highlighting}[]
\FunctionTok{cat}\NormalTok{(}\StringTok{"Desv. Est.:"}\NormalTok{, }\FunctionTok{sd}\NormalTok{(titulaciones\_agg}\SpecialCharTok{$}\NormalTok{Tasa\_Abandono), }\StringTok{"}\SpecialCharTok{\textbackslash{}n}\StringTok{"}\NormalTok{)}
\end{Highlighting}
\end{Shaded}

\begin{verbatim}
## Desv. Est.: 0.1199212
\end{verbatim}

\begin{Shaded}
\begin{Highlighting}[]
\FunctionTok{cat}\NormalTok{(}\StringTok{"Mín:"}\NormalTok{, }\FunctionTok{min}\NormalTok{(titulaciones\_agg}\SpecialCharTok{$}\NormalTok{Tasa\_Abandono), }\StringTok{"}\SpecialCharTok{\textbackslash{}n}\StringTok{"}\NormalTok{)}
\end{Highlighting}
\end{Shaded}

\begin{verbatim}
## Mín: 0
\end{verbatim}

\begin{Shaded}
\begin{Highlighting}[]
\FunctionTok{cat}\NormalTok{(}\StringTok{"Máx:"}\NormalTok{, }\FunctionTok{max}\NormalTok{(titulaciones\_agg}\SpecialCharTok{$}\NormalTok{Tasa\_Abandono), }\StringTok{"}\SpecialCharTok{\textbackslash{}n}\StringTok{"}\NormalTok{)}
\end{Highlighting}
\end{Shaded}

\begin{verbatim}
## Máx: 0.5917542
\end{verbatim}

\section{3.2 Normalización de Datos}\label{normalizaciuxf3n-de-datos}

\begin{Shaded}
\begin{Highlighting}[]
\CommentTok{\# Crear matriz de features para clustering}
\NormalTok{features\_clustering }\OtherTok{\textless{}{-}}\NormalTok{ titulaciones\_agg }\SpecialCharTok{\%\textgreater{}\%}
  \FunctionTok{select}\NormalTok{(Tasa\_Rendimiento, Tasa\_Abandono)}

\CommentTok{\# Normalizar datos (escalado Z{-}score)}
\NormalTok{features\_normalized }\OtherTok{\textless{}{-}} \FunctionTok{scale}\NormalTok{(features\_clustering)}

\CommentTok{\# Verificar normalización}
\FunctionTok{cat}\NormalTok{(}\StringTok{"Media después de normalización:}\SpecialCharTok{\textbackslash{}n}\StringTok{"}\NormalTok{)}
\end{Highlighting}
\end{Shaded}

\begin{verbatim}
## Media después de normalización:
\end{verbatim}

\begin{Shaded}
\begin{Highlighting}[]
\FunctionTok{print}\NormalTok{(}\FunctionTok{colMeans}\NormalTok{(features\_normalized))}
\end{Highlighting}
\end{Shaded}

\begin{verbatim}
## Tasa_Rendimiento    Tasa_Abandono 
##     1.580175e-16     2.395320e-17
\end{verbatim}

\begin{Shaded}
\begin{Highlighting}[]
\FunctionTok{cat}\NormalTok{(}\StringTok{"}\SpecialCharTok{\textbackslash{}n}\StringTok{Desviación estándar después de normalización:}\SpecialCharTok{\textbackslash{}n}\StringTok{"}\NormalTok{)}
\end{Highlighting}
\end{Shaded}

\begin{verbatim}
## 
## Desviación estándar después de normalización:
\end{verbatim}

\begin{Shaded}
\begin{Highlighting}[]
\FunctionTok{print}\NormalTok{(}\FunctionTok{apply}\NormalTok{(features\_normalized, }\DecValTok{2}\NormalTok{, sd))}
\end{Highlighting}
\end{Shaded}

\begin{verbatim}
## Tasa_Rendimiento    Tasa_Abandono 
##                1                1
\end{verbatim}

\begin{center}\rule{0.5\linewidth}{0.5pt}\end{center}

\section{4. Determinación del Número Óptimo de
Clústeres}\label{determinaciuxf3n-del-nuxfamero-uxf3ptimo-de-cluxfasteres}

\begin{Shaded}
\begin{Highlighting}[]
\CommentTok{\# Elbow method}
\NormalTok{wss }\OtherTok{\textless{}{-}} \FunctionTok{sapply}\NormalTok{(}\DecValTok{1}\SpecialCharTok{:}\DecValTok{10}\NormalTok{, }\ControlFlowTok{function}\NormalTok{(k) \{}
  \FunctionTok{kmeans}\NormalTok{(features\_normalized, k, }\AttributeTok{nstart =} \DecValTok{10}\NormalTok{)}\SpecialCharTok{$}\NormalTok{tot.withinss}
\NormalTok{\})}

\CommentTok{\# Gráfico Elbow}
\NormalTok{elbow\_df }\OtherTok{\textless{}{-}} \FunctionTok{data.frame}\NormalTok{(}
  \AttributeTok{k =} \DecValTok{1}\SpecialCharTok{:}\DecValTok{10}\NormalTok{,}
  \AttributeTok{wss =}\NormalTok{ wss}
\NormalTok{)}

\FunctionTok{ggplot}\NormalTok{(elbow\_df, }\FunctionTok{aes}\NormalTok{(}\AttributeTok{x =}\NormalTok{ k, }\AttributeTok{y =}\NormalTok{ wss)) }\SpecialCharTok{+}
  \FunctionTok{geom\_line}\NormalTok{(}\AttributeTok{color =} \StringTok{"\#3498db"}\NormalTok{, }\AttributeTok{size =} \DecValTok{1}\NormalTok{) }\SpecialCharTok{+}
  \FunctionTok{geom\_point}\NormalTok{(}\AttributeTok{color =} \StringTok{"\#3498db"}\NormalTok{, }\AttributeTok{size =} \DecValTok{3}\NormalTok{) }\SpecialCharTok{+}
  \FunctionTok{labs}\NormalTok{(}
    \AttributeTok{title =} \StringTok{"Método del Codo para Determinar Número Óptimo de Clústeres"}\NormalTok{,}
    \AttributeTok{x =} \StringTok{"Número de Clústeres (k)"}\NormalTok{,}
    \AttributeTok{y =} \StringTok{"Suma Intra{-}cluster de Cuadrados"}
\NormalTok{  ) }\SpecialCharTok{+}
  \FunctionTok{theme\_minimal}\NormalTok{() }\SpecialCharTok{+}
  \FunctionTok{theme}\NormalTok{(}
    \AttributeTok{plot.title =} \FunctionTok{element\_text}\NormalTok{(}\AttributeTok{hjust =} \FloatTok{0.5}\NormalTok{, }\AttributeTok{size =} \DecValTok{14}\NormalTok{, }\AttributeTok{face =} \StringTok{"bold"}\NormalTok{),}
    \AttributeTok{axis.title =} \FunctionTok{element\_text}\NormalTok{(}\AttributeTok{size =} \DecValTok{11}\NormalTok{),}
    \AttributeTok{panel.grid.major =} \FunctionTok{element\_line}\NormalTok{(}\AttributeTok{color =} \StringTok{"gray90"}\NormalTok{)}
\NormalTok{  )}
\end{Highlighting}
\end{Shaded}

\begin{center}\includegraphics{clustering_titulaciones_files/figure-latex/elbow_method-1} \end{center}

\section{4.1 Análisis de Silhueta}\label{anuxe1lisis-de-silhueta}

\begin{Shaded}
\begin{Highlighting}[]
\CommentTok{\# Análisis de silhueta para k = 2 a 6}
\NormalTok{silhouette\_scores }\OtherTok{\textless{}{-}} \FunctionTok{sapply}\NormalTok{(}\DecValTok{2}\SpecialCharTok{:}\DecValTok{6}\NormalTok{, }\ControlFlowTok{function}\NormalTok{(k) \{}
\NormalTok{  km }\OtherTok{\textless{}{-}} \FunctionTok{kmeans}\NormalTok{(features\_normalized, k, }\AttributeTok{nstart =} \DecValTok{10}\NormalTok{)}
\NormalTok{  sil }\OtherTok{\textless{}{-}} \FunctionTok{silhouette}\NormalTok{(km}\SpecialCharTok{$}\NormalTok{cluster, }\FunctionTok{dist}\NormalTok{(features\_normalized))}
  \FunctionTok{mean}\NormalTok{(sil[, }\StringTok{"sil\_width"}\NormalTok{])}
\NormalTok{\})}

\NormalTok{sil\_df }\OtherTok{\textless{}{-}} \FunctionTok{data.frame}\NormalTok{(}
  \AttributeTok{k =} \DecValTok{2}\SpecialCharTok{:}\DecValTok{6}\NormalTok{,}
  \AttributeTok{silhouette =}\NormalTok{ silhouette\_scores}
\NormalTok{)}

\FunctionTok{ggplot}\NormalTok{(sil\_df, }\FunctionTok{aes}\NormalTok{(}\AttributeTok{x =}\NormalTok{ k, }\AttributeTok{y =}\NormalTok{ silhouette)) }\SpecialCharTok{+}
  \FunctionTok{geom\_line}\NormalTok{(}\AttributeTok{color =} \StringTok{"\#e74c3c"}\NormalTok{, }\AttributeTok{size =} \DecValTok{1}\NormalTok{) }\SpecialCharTok{+}
  \FunctionTok{geom\_point}\NormalTok{(}\AttributeTok{color =} \StringTok{"\#e74c3c"}\NormalTok{, }\AttributeTok{size =} \DecValTok{3}\NormalTok{) }\SpecialCharTok{+}
  \FunctionTok{labs}\NormalTok{(}
    \AttributeTok{title =} \StringTok{"Análisis de Silhueta: Calidad de la Partición"}\NormalTok{,}
    \AttributeTok{x =} \StringTok{"Número de Clústeres (k)"}\NormalTok{,}
    \AttributeTok{y =} \StringTok{"Coeficiente de Silhueta"}
\NormalTok{  ) }\SpecialCharTok{+}
  \FunctionTok{theme\_minimal}\NormalTok{() }\SpecialCharTok{+}
  \FunctionTok{theme}\NormalTok{(}
    \AttributeTok{plot.title =} \FunctionTok{element\_text}\NormalTok{(}\AttributeTok{hjust =} \FloatTok{0.5}\NormalTok{, }\AttributeTok{size =} \DecValTok{14}\NormalTok{, }\AttributeTok{face =} \StringTok{"bold"}\NormalTok{),}
    \AttributeTok{axis.title =} \FunctionTok{element\_text}\NormalTok{(}\AttributeTok{size =} \DecValTok{11}\NormalTok{),}
    \AttributeTok{panel.grid.major =} \FunctionTok{element\_line}\NormalTok{(}\AttributeTok{color =} \StringTok{"gray90"}\NormalTok{)}
\NormalTok{  )}
\end{Highlighting}
\end{Shaded}

\begin{center}\includegraphics{clustering_titulaciones_files/figure-latex/silhouette_analysis-1} \end{center}

\begin{Shaded}
\begin{Highlighting}[]
\FunctionTok{cat}\NormalTok{(}\StringTok{"Scores de Silhueta:}\SpecialCharTok{\textbackslash{}n}\StringTok{"}\NormalTok{)}
\end{Highlighting}
\end{Shaded}

\begin{verbatim}
## Scores de Silhueta:
\end{verbatim}

\begin{Shaded}
\begin{Highlighting}[]
\FunctionTok{print}\NormalTok{(sil\_df)}
\end{Highlighting}
\end{Shaded}

\begin{verbatim}
##   k silhouette
## 1 2  0.5382338
## 2 3  0.4444979
## 3 4  0.3223800
## 4 5  0.3375820
## 5 6  0.3230141
\end{verbatim}

\begin{Shaded}
\begin{Highlighting}[]
\FunctionTok{cat}\NormalTok{(}\StringTok{"}\SpecialCharTok{\textbackslash{}n}\StringTok{"}\NormalTok{)}
\end{Highlighting}
\end{Shaded}

\begin{Shaded}
\begin{Highlighting}[]
\FunctionTok{cat}\NormalTok{(}\StringTok{"Silhouette Scores por k:}\SpecialCharTok{\textbackslash{}n}\StringTok{"}\NormalTok{)}
\end{Highlighting}
\end{Shaded}

\begin{verbatim}
## Silhouette Scores por k:
\end{verbatim}

\begin{Shaded}
\begin{Highlighting}[]
\FunctionTok{print}\NormalTok{(sil\_df)}
\end{Highlighting}
\end{Shaded}

\begin{verbatim}
##   k silhouette
## 1 2  0.5382338
## 2 3  0.4444979
## 3 4  0.3223800
## 4 5  0.3375820
## 5 6  0.3230141
\end{verbatim}

\begin{Shaded}
\begin{Highlighting}[]
\FunctionTok{cat}\NormalTok{(}\StringTok{"}\SpecialCharTok{\textbackslash{}n}\StringTok{k óptimo según Silhouette:"}\NormalTok{, sil\_df}\SpecialCharTok{$}\NormalTok{k[}\FunctionTok{which.max}\NormalTok{(sil\_df}\SpecialCharTok{$}\NormalTok{silhouette)], }\StringTok{"}\SpecialCharTok{\textbackslash{}n}\StringTok{"}\NormalTok{)}
\end{Highlighting}
\end{Shaded}

\begin{verbatim}
## 
## k óptimo según Silhouette: 2
\end{verbatim}

\begin{center}\rule{0.5\linewidth}{0.5pt}\end{center}

\section{5. Aplicación de K-means
Clustering}\label{aplicaciuxf3n-de-k-means-clustering}

\begin{Shaded}
\begin{Highlighting}[]
\CommentTok{\# Determinar k óptimo basado en silhouette analysis}
\NormalTok{silhouette\_scores\_optimal }\OtherTok{\textless{}{-}} \FunctionTok{sapply}\NormalTok{(}\DecValTok{2}\SpecialCharTok{:}\DecValTok{6}\NormalTok{, }\ControlFlowTok{function}\NormalTok{(k) \{}
\NormalTok{  km }\OtherTok{\textless{}{-}} \FunctionTok{kmeans}\NormalTok{(features\_normalized, k, }\AttributeTok{nstart =} \DecValTok{10}\NormalTok{)}
\NormalTok{  sil }\OtherTok{\textless{}{-}} \FunctionTok{silhouette}\NormalTok{(km}\SpecialCharTok{$}\NormalTok{cluster, }\FunctionTok{dist}\NormalTok{(features\_normalized))}
  \FunctionTok{mean}\NormalTok{(sil[, }\StringTok{"sil\_width"}\NormalTok{])}
\NormalTok{\})}

\NormalTok{k\_optimal\_value }\OtherTok{\textless{}{-}} \FunctionTok{which.max}\NormalTok{(silhouette\_scores\_optimal) }\SpecialCharTok{+} \DecValTok{1}  \CommentTok{\# +1 porque comienza en k=2}
\FunctionTok{cat}\NormalTok{(}\StringTok{"K ÓPTIMO DETERMINADO:"}\NormalTok{, k\_optimal\_value, }\StringTok{"}\SpecialCharTok{\textbackslash{}n}\StringTok{"}\NormalTok{)}
\end{Highlighting}
\end{Shaded}

\begin{verbatim}
## K ÓPTIMO DETERMINADO: 2
\end{verbatim}

\begin{Shaded}
\begin{Highlighting}[]
\FunctionTok{cat}\NormalTok{(}\StringTok{"Silhouette Score:"}\NormalTok{, }\FunctionTok{round}\NormalTok{(}\FunctionTok{max}\NormalTok{(silhouette\_scores\_optimal), }\DecValTok{4}\NormalTok{), }\StringTok{"}\SpecialCharTok{\textbackslash{}n\textbackslash{}n}\StringTok{"}\NormalTok{)}
\end{Highlighting}
\end{Shaded}

\begin{verbatim}
## Silhouette Score: 0.5382
\end{verbatim}

\begin{Shaded}
\begin{Highlighting}[]
\CommentTok{\# Usar k óptimo para clustering final}
\FunctionTok{set.seed}\NormalTok{(}\DecValTok{123}\NormalTok{)}
\NormalTok{k\_optimal }\OtherTok{\textless{}{-}}\NormalTok{ k\_optimal\_value}
\NormalTok{kmeans\_result }\OtherTok{\textless{}{-}} \FunctionTok{kmeans}\NormalTok{(features\_normalized, }\AttributeTok{centers =}\NormalTok{ k\_optimal, }\AttributeTok{nstart =} \DecValTok{25}\NormalTok{, }\AttributeTok{iter.max =} \DecValTok{300}\NormalTok{)}

\CommentTok{\# Crear etiquetas dinámicas para clusters}
\NormalTok{cluster\_labels }\OtherTok{\textless{}{-}} \FunctionTok{paste0}\NormalTok{(}\StringTok{"Clúster "}\NormalTok{, }\DecValTok{1}\SpecialCharTok{:}\NormalTok{k\_optimal)}

\CommentTok{\# Añadir clústeres a los datos originales}
\NormalTok{titulaciones\_clustered }\OtherTok{\textless{}{-}}\NormalTok{ titulaciones\_agg }\SpecialCharTok{\%\textgreater{}\%}
  \FunctionTok{mutate}\NormalTok{(}
    \AttributeTok{Cluster =} \FunctionTok{as.factor}\NormalTok{(kmeans\_result}\SpecialCharTok{$}\NormalTok{cluster),}
    \AttributeTok{Cluster\_Label =} \FunctionTok{paste0}\NormalTok{(}\StringTok{"Clúster "}\NormalTok{, Cluster)}
\NormalTok{  )}

\CommentTok{\# Resumen de clústeres}
\FunctionTok{cat}\NormalTok{(}\StringTok{"RESUMEN DE CLÚSTERES}\SpecialCharTok{\textbackslash{}n}\StringTok{"}\NormalTok{)}
\end{Highlighting}
\end{Shaded}

\begin{verbatim}
## RESUMEN DE CLÚSTERES
\end{verbatim}

\begin{Shaded}
\begin{Highlighting}[]
\FunctionTok{cat}\NormalTok{(}\StringTok{"Número de titulaciones por clúster:}\SpecialCharTok{\textbackslash{}n}\StringTok{"}\NormalTok{)}
\end{Highlighting}
\end{Shaded}

\begin{verbatim}
## Número de titulaciones por clúster:
\end{verbatim}

\begin{Shaded}
\begin{Highlighting}[]
\FunctionTok{print}\NormalTok{(}\FunctionTok{table}\NormalTok{(titulaciones\_clustered}\SpecialCharTok{$}\NormalTok{Cluster))}
\end{Highlighting}
\end{Shaded}

\begin{verbatim}
## 
##   1   2 
##  59 139
\end{verbatim}

\section{5.1 Características de Cada
Clúster}\label{caracteruxedsticas-de-cada-cluxfaster}

\begin{Shaded}
\begin{Highlighting}[]
\NormalTok{cluster\_summary }\OtherTok{\textless{}{-}}\NormalTok{ titulaciones\_clustered }\SpecialCharTok{\%\textgreater{}\%}
  \FunctionTok{group\_by}\NormalTok{(Cluster) }\SpecialCharTok{\%\textgreater{}\%}
  \FunctionTok{summarise}\NormalTok{(}
    \AttributeTok{n\_titulaciones =} \FunctionTok{n}\NormalTok{(),}
    \AttributeTok{Rendimiento\_Media =} \FunctionTok{mean}\NormalTok{(Tasa\_Rendimiento),}
    \AttributeTok{Rendimiento\_SD =} \FunctionTok{sd}\NormalTok{(Tasa\_Rendimiento),}
    \AttributeTok{Abandono\_Media =} \FunctionTok{mean}\NormalTok{(Tasa\_Abandono),}
    \AttributeTok{Abandono\_SD =} \FunctionTok{sd}\NormalTok{(Tasa\_Abandono),}
    \AttributeTok{.groups =} \StringTok{\textquotesingle{}drop\textquotesingle{}}
\NormalTok{  )}

\FunctionTok{print}\NormalTok{(cluster\_summary)}
\end{Highlighting}
\end{Shaded}

\begin{verbatim}
## # A tibble: 2 x 6
##   Cluster n_titulaciones Rendimiento_Media Rendimiento_SD Abandono_Media
##   <fct>            <int>             <dbl>          <dbl>          <dbl>
## 1 1                   59             0.593         0.0866          0.311
## 2 2                  139             0.811         0.0854          0.107
## # i 1 more variable: Abandono_SD <dbl>
\end{verbatim}

\begin{Shaded}
\begin{Highlighting}[]
\CommentTok{\# Tabla detallada con mejor formato}
\NormalTok{knitr}\SpecialCharTok{::}\FunctionTok{kable}\NormalTok{(cluster\_summary, }\AttributeTok{digits =} \DecValTok{4}\NormalTok{, }\AttributeTok{caption =} \StringTok{"Características de Clústeres"}\NormalTok{)}
\end{Highlighting}
\end{Shaded}

\begin{longtable}[]{@{}
  >{\raggedright\arraybackslash}p{(\linewidth - 10\tabcolsep) * \real{0.0964}}
  >{\raggedleft\arraybackslash}p{(\linewidth - 10\tabcolsep) * \real{0.1807}}
  >{\raggedleft\arraybackslash}p{(\linewidth - 10\tabcolsep) * \real{0.2169}}
  >{\raggedleft\arraybackslash}p{(\linewidth - 10\tabcolsep) * \real{0.1807}}
  >{\raggedleft\arraybackslash}p{(\linewidth - 10\tabcolsep) * \real{0.1807}}
  >{\raggedleft\arraybackslash}p{(\linewidth - 10\tabcolsep) * \real{0.1446}}@{}}
\caption{Características de Clústeres}\tabularnewline
\toprule\noalign{}
\begin{minipage}[b]{\linewidth}\raggedright
Cluster
\end{minipage} & \begin{minipage}[b]{\linewidth}\raggedleft
n\_titulaciones
\end{minipage} & \begin{minipage}[b]{\linewidth}\raggedleft
Rendimiento\_Media
\end{minipage} & \begin{minipage}[b]{\linewidth}\raggedleft
Rendimiento\_SD
\end{minipage} & \begin{minipage}[b]{\linewidth}\raggedleft
Abandono\_Media
\end{minipage} & \begin{minipage}[b]{\linewidth}\raggedleft
Abandono\_SD
\end{minipage} \\
\midrule\noalign{}
\endfirsthead
\toprule\noalign{}
\begin{minipage}[b]{\linewidth}\raggedright
Cluster
\end{minipage} & \begin{minipage}[b]{\linewidth}\raggedleft
n\_titulaciones
\end{minipage} & \begin{minipage}[b]{\linewidth}\raggedleft
Rendimiento\_Media
\end{minipage} & \begin{minipage}[b]{\linewidth}\raggedleft
Rendimiento\_SD
\end{minipage} & \begin{minipage}[b]{\linewidth}\raggedleft
Abandono\_Media
\end{minipage} & \begin{minipage}[b]{\linewidth}\raggedleft
Abandono\_SD
\end{minipage} \\
\midrule\noalign{}
\endhead
\bottomrule\noalign{}
\endlastfoot
1 & 59 & 0.5927 & 0.0866 & 0.3106 & 0.1084 \\
2 & 139 & 0.8108 & 0.0854 & 0.1071 & 0.0562 \\
\end{longtable}

\begin{center}\rule{0.5\linewidth}{0.5pt}\end{center}

\section{6. Visualización Principal: Cluster
Plot}\label{visualizaciuxf3n-principal-cluster-plot}

\begin{Shaded}
\begin{Highlighting}[]
\CommentTok{\# Crear paleta de colores dinámicamente según k\_optimal}
\NormalTok{paleta\_colores }\OtherTok{\textless{}{-}} \FunctionTok{c}\NormalTok{(}\StringTok{"\#e74c3c"}\NormalTok{, }\StringTok{"\#3498db"}\NormalTok{, }\StringTok{"\#2ecc71"}\NormalTok{, }\StringTok{"\#f39c12"}\NormalTok{, }\StringTok{"\#9b59b6"}\NormalTok{, }\StringTok{"\#1abc9c"}\NormalTok{)}
\NormalTok{colores }\OtherTok{\textless{}{-}}\NormalTok{ paleta\_colores[}\DecValTok{1}\SpecialCharTok{:}\NormalTok{k\_optimal]}

\CommentTok{\# Gráfico principal de clústeres}
\NormalTok{plot\_clusters }\OtherTok{\textless{}{-}} \FunctionTok{ggplot}\NormalTok{(titulaciones\_clustered, }\FunctionTok{aes}\NormalTok{(}
  \AttributeTok{x =}\NormalTok{ Tasa\_Rendimiento,}
  \AttributeTok{y =}\NormalTok{ Tasa\_Abandono,}
  \AttributeTok{color =}\NormalTok{ Cluster,}
  \AttributeTok{label =} \FunctionTok{str\_wrap}\NormalTok{(Titulacion, }\DecValTok{20}\NormalTok{)}
\NormalTok{)) }\SpecialCharTok{+}
  \FunctionTok{geom\_point}\NormalTok{(}\AttributeTok{size =} \DecValTok{6}\NormalTok{, }\AttributeTok{alpha =} \FloatTok{0.7}\NormalTok{) }\SpecialCharTok{+}
  \FunctionTok{geom\_label\_repel}\NormalTok{(}
    \AttributeTok{size =} \DecValTok{3}\NormalTok{,}
    \AttributeTok{max.overlaps =} \DecValTok{30}\NormalTok{,}
    \AttributeTok{box.padding =} \FloatTok{0.5}\NormalTok{,}
    \AttributeTok{point.padding =} \FloatTok{0.5}
\NormalTok{  ) }\SpecialCharTok{+}
  \CommentTok{\# Añadir centroides}
  \FunctionTok{geom\_point}\NormalTok{(}
    \AttributeTok{data =} \FunctionTok{as.data.frame}\NormalTok{(kmeans\_result}\SpecialCharTok{$}\NormalTok{centers) }\SpecialCharTok{\%\textgreater{}\%}
      \FunctionTok{mutate}\NormalTok{(}
        \AttributeTok{Tasa\_Rendimiento =}\NormalTok{ Tasa\_Rendimiento,}
        \AttributeTok{Tasa\_Abandono =}\NormalTok{ Tasa\_Abandono,}
        \AttributeTok{Cluster =} \FunctionTok{factor}\NormalTok{(}\DecValTok{1}\SpecialCharTok{:}\NormalTok{k\_optimal)}
\NormalTok{      ),}
    \FunctionTok{aes}\NormalTok{(}\AttributeTok{x =}\NormalTok{ Tasa\_Rendimiento, }\AttributeTok{y =}\NormalTok{ Tasa\_Abandono, }\AttributeTok{color =}\NormalTok{ Cluster),}
    \AttributeTok{shape =} \DecValTok{17}\NormalTok{,}
    \AttributeTok{size =} \DecValTok{10}\NormalTok{,}
    \AttributeTok{inherit.aes =} \ConstantTok{FALSE}
\NormalTok{  ) }\SpecialCharTok{+}
  \FunctionTok{scale\_color\_manual}\NormalTok{(}\AttributeTok{values =}\NormalTok{ colores, }\AttributeTok{labels =}\NormalTok{ cluster\_labels) }\SpecialCharTok{+}
  \FunctionTok{labs}\NormalTok{(}
    \AttributeTok{title =} \StringTok{"Agrupamiento de Titulaciones según Rendimiento y Abandono"}\NormalTok{,}
    \AttributeTok{subtitle =} \StringTok{"Triángulos representan los centroides de cada clúster"}\NormalTok{,}
    \AttributeTok{x =} \StringTok{"Tasa de Rendimiento (media del período)"}\NormalTok{,}
    \AttributeTok{y =} \StringTok{"Tasa de Abandono (media del período)"}\NormalTok{,}
    \AttributeTok{color =} \StringTok{"Clúster"}
\NormalTok{  ) }\SpecialCharTok{+}
  \FunctionTok{theme\_minimal}\NormalTok{() }\SpecialCharTok{+}
  \FunctionTok{theme}\NormalTok{(}
    \AttributeTok{plot.title =} \FunctionTok{element\_text}\NormalTok{(}\AttributeTok{hjust =} \FloatTok{0.5}\NormalTok{, }\AttributeTok{size =} \DecValTok{16}\NormalTok{, }\AttributeTok{face =} \StringTok{"bold"}\NormalTok{),}
    \AttributeTok{plot.subtitle =} \FunctionTok{element\_text}\NormalTok{(}\AttributeTok{hjust =} \FloatTok{0.5}\NormalTok{, }\AttributeTok{size =} \DecValTok{12}\NormalTok{, }\AttributeTok{face =} \StringTok{"italic"}\NormalTok{),}
    \AttributeTok{axis.title =} \FunctionTok{element\_text}\NormalTok{(}\AttributeTok{size =} \DecValTok{12}\NormalTok{, }\AttributeTok{face =} \StringTok{"bold"}\NormalTok{),}
    \AttributeTok{legend.position =} \StringTok{"right"}\NormalTok{,}
    \AttributeTok{panel.grid.major =} \FunctionTok{element\_line}\NormalTok{(}\AttributeTok{color =} \StringTok{"gray90"}\NormalTok{),}
    \AttributeTok{panel.grid.minor =} \FunctionTok{element\_blank}\NormalTok{()}
\NormalTok{  )}

\FunctionTok{print}\NormalTok{(plot\_clusters)}
\end{Highlighting}
\end{Shaded}

\begin{center}\includegraphics{clustering_titulaciones_files/figure-latex/cluster_plot-1} \end{center}

\section{6.1 Gráfico Alternativo sin
Etiquetas}\label{gruxe1fico-alternativo-sin-etiquetas}

\begin{Shaded}
\begin{Highlighting}[]
\CommentTok{\# Calcular medianas para dividir en cuadrantes}
\NormalTok{median\_rend }\OtherTok{\textless{}{-}} \FunctionTok{median}\NormalTok{(titulaciones\_clustered}\SpecialCharTok{$}\NormalTok{Tasa\_Rendimiento)}
\NormalTok{median\_aband }\OtherTok{\textless{}{-}} \FunctionTok{median}\NormalTok{(titulaciones\_clustered}\SpecialCharTok{$}\NormalTok{Tasa\_Abandono)}

\CommentTok{\# Versión con 4 cuadrantes claramente definidos}
\NormalTok{plot\_clusters\_clean }\OtherTok{\textless{}{-}} \FunctionTok{ggplot}\NormalTok{(titulaciones\_clustered, }\FunctionTok{aes}\NormalTok{(}
  \AttributeTok{x =}\NormalTok{ Tasa\_Rendimiento,}
  \AttributeTok{y =}\NormalTok{ Tasa\_Abandono,}
  \AttributeTok{color =}\NormalTok{ Cluster,}
  \AttributeTok{fill =}\NormalTok{ Cluster}
\NormalTok{)) }\SpecialCharTok{+}
  \CommentTok{\# Fondo de cuadrantes con transparencia}
  \FunctionTok{geom\_rect}\NormalTok{(}
    \FunctionTok{aes}\NormalTok{(}\AttributeTok{xmin =} \SpecialCharTok{{-}}\ConstantTok{Inf}\NormalTok{, }\AttributeTok{xmax =}\NormalTok{ median\_rend, }\AttributeTok{ymin =} \SpecialCharTok{{-}}\ConstantTok{Inf}\NormalTok{, }\AttributeTok{ymax =}\NormalTok{ median\_aband),}
    \AttributeTok{fill =} \StringTok{"\#ecf0f1"}\NormalTok{, }\AttributeTok{alpha =} \FloatTok{0.3}\NormalTok{, }\AttributeTok{inherit.aes =} \ConstantTok{FALSE}
\NormalTok{  ) }\SpecialCharTok{+}
  \FunctionTok{geom\_rect}\NormalTok{(}
    \FunctionTok{aes}\NormalTok{(}\AttributeTok{xmin =}\NormalTok{ median\_rend, }\AttributeTok{xmax =} \ConstantTok{Inf}\NormalTok{, }\AttributeTok{ymin =} \SpecialCharTok{{-}}\ConstantTok{Inf}\NormalTok{, }\AttributeTok{ymax =}\NormalTok{ median\_aband),}
    \AttributeTok{fill =} \StringTok{"\#d5f4e6"}\NormalTok{, }\AttributeTok{alpha =} \FloatTok{0.3}\NormalTok{, }\AttributeTok{inherit.aes =} \ConstantTok{FALSE}
\NormalTok{  ) }\SpecialCharTok{+}
  \FunctionTok{geom\_rect}\NormalTok{(}
    \FunctionTok{aes}\NormalTok{(}\AttributeTok{xmin =} \SpecialCharTok{{-}}\ConstantTok{Inf}\NormalTok{, }\AttributeTok{xmax =}\NormalTok{ median\_rend, }\AttributeTok{ymin =}\NormalTok{ median\_aband, }\AttributeTok{ymax =} \ConstantTok{Inf}\NormalTok{),}
    \AttributeTok{fill =} \StringTok{"\#fadbd8"}\NormalTok{, }\AttributeTok{alpha =} \FloatTok{0.3}\NormalTok{, }\AttributeTok{inherit.aes =} \ConstantTok{FALSE}
\NormalTok{  ) }\SpecialCharTok{+}
  \FunctionTok{geom\_rect}\NormalTok{(}
    \FunctionTok{aes}\NormalTok{(}\AttributeTok{xmin =}\NormalTok{ median\_rend, }\AttributeTok{xmax =} \ConstantTok{Inf}\NormalTok{, }\AttributeTok{ymin =}\NormalTok{ median\_aband, }\AttributeTok{ymax =} \ConstantTok{Inf}\NormalTok{),}
    \AttributeTok{fill =} \StringTok{"\#d6eaf8"}\NormalTok{, }\AttributeTok{alpha =} \FloatTok{0.3}\NormalTok{, }\AttributeTok{inherit.aes =} \ConstantTok{FALSE}
\NormalTok{  ) }\SpecialCharTok{+}
  \CommentTok{\# Líneas de división}
  \FunctionTok{geom\_vline}\NormalTok{(}\AttributeTok{xintercept =}\NormalTok{ median\_rend, }\AttributeTok{linetype =} \StringTok{"dashed"}\NormalTok{, }\AttributeTok{color =} \StringTok{"gray40"}\NormalTok{, }\AttributeTok{alpha =} \FloatTok{0.7}\NormalTok{, }\AttributeTok{size =} \DecValTok{1}\NormalTok{) }\SpecialCharTok{+}
  \FunctionTok{geom\_hline}\NormalTok{(}\AttributeTok{yintercept =}\NormalTok{ median\_aband, }\AttributeTok{linetype =} \StringTok{"dashed"}\NormalTok{, }\AttributeTok{color =} \StringTok{"gray40"}\NormalTok{, }\AttributeTok{alpha =} \FloatTok{0.7}\NormalTok{, }\AttributeTok{size =} \DecValTok{1}\NormalTok{) }\SpecialCharTok{+}
  \CommentTok{\# Puntos}
  \FunctionTok{geom\_point}\NormalTok{(}\AttributeTok{size =} \DecValTok{8}\NormalTok{, }\AttributeTok{alpha =} \FloatTok{0.8}\NormalTok{, }\AttributeTok{shape =} \DecValTok{21}\NormalTok{, }\AttributeTok{color =} \StringTok{"white"}\NormalTok{, }\AttributeTok{stroke =} \DecValTok{2}\NormalTok{) }\SpecialCharTok{+}
  \CommentTok{\# Etiquetas de cuadrantes}
  \FunctionTok{annotate}\NormalTok{(}
    \StringTok{"text"}\NormalTok{,}
    \AttributeTok{x =} \FunctionTok{min}\NormalTok{(titulaciones\_clustered}\SpecialCharTok{$}\NormalTok{Tasa\_Rendimiento) }\SpecialCharTok{+} \FloatTok{0.02}\NormalTok{,}
    \AttributeTok{y =} \FunctionTok{max}\NormalTok{(titulaciones\_clustered}\SpecialCharTok{$}\NormalTok{Tasa\_Abandono) }\SpecialCharTok{{-}} \FloatTok{0.02}\NormalTok{,}
    \AttributeTok{label =} \StringTok{"BAJO RENDIMIENTO}\SpecialCharTok{\textbackslash{}n}\StringTok{ALTO ABANDONO}\SpecialCharTok{\textbackslash{}n}\StringTok{(EN RIESGO)"}\NormalTok{,}
    \AttributeTok{hjust =} \DecValTok{0}\NormalTok{, }\AttributeTok{vjust =} \DecValTok{1}\NormalTok{,}
    \AttributeTok{size =} \DecValTok{4}\NormalTok{,}
    \AttributeTok{color =} \StringTok{"\#c0392b"}\NormalTok{,}
    \AttributeTok{fontface =} \StringTok{"bold"}\NormalTok{,}
    \AttributeTok{alpha =} \FloatTok{0.8}
\NormalTok{  ) }\SpecialCharTok{+}
  \FunctionTok{annotate}\NormalTok{(}
    \StringTok{"text"}\NormalTok{,}
    \AttributeTok{x =} \FunctionTok{max}\NormalTok{(titulaciones\_clustered}\SpecialCharTok{$}\NormalTok{Tasa\_Rendimiento) }\SpecialCharTok{{-}} \FloatTok{0.02}\NormalTok{,}
    \AttributeTok{y =} \FunctionTok{max}\NormalTok{(titulaciones\_clustered}\SpecialCharTok{$}\NormalTok{Tasa\_Abandono) }\SpecialCharTok{{-}} \FloatTok{0.02}\NormalTok{,}
    \AttributeTok{label =} \StringTok{"ALTO RENDIMIENTO}\SpecialCharTok{\textbackslash{}n}\StringTok{ALTO ABANDONO}\SpecialCharTok{\textbackslash{}n}\StringTok{(ABANDONO SELECTIVO)"}\NormalTok{,}
    \AttributeTok{hjust =} \DecValTok{1}\NormalTok{, }\AttributeTok{vjust =} \DecValTok{1}\NormalTok{,}
    \AttributeTok{size =} \DecValTok{4}\NormalTok{,}
    \AttributeTok{color =} \StringTok{"\#8e44ad"}\NormalTok{,}
    \AttributeTok{fontface =} \StringTok{"bold"}\NormalTok{,}
    \AttributeTok{alpha =} \FloatTok{0.8}
\NormalTok{  ) }\SpecialCharTok{+}
  \FunctionTok{annotate}\NormalTok{(}
    \StringTok{"text"}\NormalTok{,}
    \AttributeTok{x =} \FunctionTok{min}\NormalTok{(titulaciones\_clustered}\SpecialCharTok{$}\NormalTok{Tasa\_Rendimiento) }\SpecialCharTok{+} \FloatTok{0.02}\NormalTok{,}
    \AttributeTok{y =} \FunctionTok{min}\NormalTok{(titulaciones\_clustered}\SpecialCharTok{$}\NormalTok{Tasa\_Abandono) }\SpecialCharTok{+} \FloatTok{0.02}\NormalTok{,}
    \AttributeTok{label =} \StringTok{"BAJO RENDIMIENTO}\SpecialCharTok{\textbackslash{}n}\StringTok{BAJO ABANDONO}\SpecialCharTok{\textbackslash{}n}\StringTok{(PROBLEMÁTICO)"}\NormalTok{,}
    \AttributeTok{hjust =} \DecValTok{0}\NormalTok{, }\AttributeTok{vjust =} \DecValTok{0}\NormalTok{,}
    \AttributeTok{size =} \DecValTok{4}\NormalTok{,}
    \AttributeTok{color =} \StringTok{"\#e67e22"}\NormalTok{,}
    \AttributeTok{fontface =} \StringTok{"bold"}\NormalTok{,}
    \AttributeTok{alpha =} \FloatTok{0.8}
\NormalTok{  ) }\SpecialCharTok{+}
  \FunctionTok{annotate}\NormalTok{(}
    \StringTok{"text"}\NormalTok{,}
    \AttributeTok{x =} \FunctionTok{max}\NormalTok{(titulaciones\_clustered}\SpecialCharTok{$}\NormalTok{Tasa\_Rendimiento) }\SpecialCharTok{{-}} \FloatTok{0.02}\NormalTok{,}
    \AttributeTok{y =} \FunctionTok{min}\NormalTok{(titulaciones\_clustered}\SpecialCharTok{$}\NormalTok{Tasa\_Abandono) }\SpecialCharTok{+} \FloatTok{0.02}\NormalTok{,}
    \AttributeTok{label =} \StringTok{"ALTO RENDIMIENTO}\SpecialCharTok{\textbackslash{}n}\StringTok{BAJO ABANDONO}\SpecialCharTok{\textbackslash{}n}\StringTok{(EXITOSAS)"}\NormalTok{,}
    \AttributeTok{hjust =} \DecValTok{1}\NormalTok{, }\AttributeTok{vjust =} \DecValTok{0}\NormalTok{,}
    \AttributeTok{size =} \DecValTok{4}\NormalTok{,}
    \AttributeTok{color =} \StringTok{"\#27ae60"}\NormalTok{,}
    \AttributeTok{fontface =} \StringTok{"bold"}\NormalTok{,}
    \AttributeTok{alpha =} \FloatTok{0.8}
\NormalTok{  ) }\SpecialCharTok{+}
  \FunctionTok{scale\_color\_manual}\NormalTok{(}\AttributeTok{values =}\NormalTok{ colores) }\SpecialCharTok{+}
  \FunctionTok{scale\_fill\_manual}\NormalTok{(}\AttributeTok{values =}\NormalTok{ colores) }\SpecialCharTok{+}
  \FunctionTok{scale\_x\_continuous}\NormalTok{(}\AttributeTok{labels =}\NormalTok{ scales}\SpecialCharTok{::}\NormalTok{percent) }\SpecialCharTok{+}
  \FunctionTok{scale\_y\_continuous}\NormalTok{(}\AttributeTok{labels =}\NormalTok{ scales}\SpecialCharTok{::}\NormalTok{percent) }\SpecialCharTok{+}
  \FunctionTok{labs}\NormalTok{(}
    \AttributeTok{title =} \StringTok{"Matriz de Posicionamiento de Titulaciones: 4 Cuadrantes"}\NormalTok{,}
    \AttributeTok{subtitle =} \StringTok{"Clasificación según Tasa de Rendimiento y Abandono (medianas como punto de corte)"}\NormalTok{,}
    \AttributeTok{x =} \StringTok{"Tasa de Rendimiento"}\NormalTok{,}
    \AttributeTok{y =} \StringTok{"Tasa de Abandono"}\NormalTok{,}
    \AttributeTok{color =} \StringTok{"Clúster"}\NormalTok{,}
    \AttributeTok{fill =} \StringTok{"Clúster"}
\NormalTok{  ) }\SpecialCharTok{+}
  \FunctionTok{theme\_minimal}\NormalTok{() }\SpecialCharTok{+}
  \FunctionTok{theme}\NormalTok{(}
    \AttributeTok{plot.title =} \FunctionTok{element\_text}\NormalTok{(}\AttributeTok{hjust =} \FloatTok{0.5}\NormalTok{, }\AttributeTok{size =} \DecValTok{16}\NormalTok{, }\AttributeTok{face =} \StringTok{"bold"}\NormalTok{),}
    \AttributeTok{plot.subtitle =} \FunctionTok{element\_text}\NormalTok{(}\AttributeTok{hjust =} \FloatTok{0.5}\NormalTok{, }\AttributeTok{size =} \DecValTok{11}\NormalTok{, }\AttributeTok{face =} \StringTok{"italic"}\NormalTok{),}
    \AttributeTok{axis.title =} \FunctionTok{element\_text}\NormalTok{(}\AttributeTok{size =} \DecValTok{12}\NormalTok{, }\AttributeTok{face =} \StringTok{"bold"}\NormalTok{),}
    \AttributeTok{axis.text =} \FunctionTok{element\_text}\NormalTok{(}\AttributeTok{size =} \DecValTok{11}\NormalTok{),}
    \AttributeTok{legend.position =} \StringTok{"right"}\NormalTok{,}
    \AttributeTok{panel.grid.major =} \FunctionTok{element\_line}\NormalTok{(}\AttributeTok{color =} \StringTok{"gray90"}\NormalTok{, }\AttributeTok{size =} \FloatTok{0.3}\NormalTok{),}
    \AttributeTok{panel.grid.minor =} \FunctionTok{element\_blank}\NormalTok{(),}
    \AttributeTok{plot.background =} \FunctionTok{element\_rect}\NormalTok{(}\AttributeTok{fill =} \StringTok{"white"}\NormalTok{, }\AttributeTok{color =} \ConstantTok{NA}\NormalTok{)}
\NormalTok{  )}

\FunctionTok{print}\NormalTok{(plot\_clusters\_clean)}
\end{Highlighting}
\end{Shaded}

\begin{center}\includegraphics{clustering_titulaciones_files/figure-latex/cluster_plot_clean-1} \end{center}

\begin{Shaded}
\begin{Highlighting}[]
\CommentTok{\# Mostrar las medianas}
\FunctionTok{cat}\NormalTok{(}\StringTok{"Mediana de Tasa de Rendimiento:"}\NormalTok{, scales}\SpecialCharTok{::}\FunctionTok{percent}\NormalTok{(median\_rend), }\StringTok{"}\SpecialCharTok{\textbackslash{}n}\StringTok{"}\NormalTok{)}
\end{Highlighting}
\end{Shaded}

\begin{verbatim}
## Mediana de Tasa de Rendimiento: 76%
\end{verbatim}

\begin{Shaded}
\begin{Highlighting}[]
\FunctionTok{cat}\NormalTok{(}\StringTok{"Mediana de Tasa de Abandono:"}\NormalTok{, scales}\SpecialCharTok{::}\FunctionTok{percent}\NormalTok{(median\_aband), }\StringTok{"}\SpecialCharTok{\textbackslash{}n}\StringTok{"}\NormalTok{)}
\end{Highlighting}
\end{Shaded}

\begin{verbatim}
## Mediana de Tasa de Abandono: 13%
\end{verbatim}

\begin{center}\rule{0.5\linewidth}{0.5pt}\end{center}

\section{6.2 Aclaración: ¿Por qué 2 Clusters pero 4
Cuadrantes?}\label{aclaraciuxf3n-por-quuxe9-2-clusters-pero-4-cuadrantes}

Una pregunta importante surge al observar que el algoritmo K-means
frecuentemente identifica \textbf{2 clusters óptimos}, mientras que
nuestra visualización propone \textbf{4 cuadrantes}. Esto es
completamente normal y refleja enfoques complementarios:

\subsection{¿Qué significa si K-means elige 2
clusters?}\label{quuxe9-significa-si-k-means-elige-2-clusters}

Si el análisis de silhueta determina que \textbf{k=2 es óptimo}, esto
generalmente indica:

\begin{enumerate}
\def\labelenumi{\arabic{enumi}.}
\item
  \textbf{Dicotomía Principal}: Existe una división fundamental en las
  titulaciones:

  \begin{itemize}
  \tightlist
  \item
    \textbf{Cluster 1}: Titulaciones exitosas (alto rendimiento, bajo
    abandono)
  \item
    \textbf{Cluster 2}: Titulaciones problemáticas (bajo rendimiento,
    alto abandono)
  \end{itemize}
\item
  \textbf{Distribución de los Datos}: La mayoría de titulaciones tienden
  a concentrarse en dos polos, no hay un ``grupo intermedio''
  estadísticamente significativo.
\item
  \textbf{Varianza Explicada}: Dos grupos explican la mayoría de la
  variabilidad, reduciendo la complejidad innecesaria.
\end{enumerate}

\subsection{¿Por qué usamos 4 cuadrantes
entonces?}\label{por-quuxe9-usamos-4-cuadrantes-entonces}

Los 4 cuadrantes, aunque no sean confirmados por K-means como clusters
autónomos, proporcionan \textbf{información estratégica adicional}:

\begin{itemize}
\item
  \textbf{Casos matizados}: Las titulaciones en el cuadrante ``Abandono
  Selectivo'' (Q2) tienen alto rendimiento pero alto abandono -
  requieren intervención \emph{diferente} que las en riesgo (Q4).
\item
  \textbf{Plan de acción operativo}: Para un gestor académico, los 4
  cuadrantes ofrecen una matriz clara y actionable.
\item
  \textbf{Perspectiva complementaria}: Se pueden interpretar como:

  \begin{itemize}
  \tightlist
  \item
    \textbf{K-means}: ¿Qué patrones naturales hay en los datos?
  \item
    \textbf{Cuadrantes}: ¿Cómo debemos intervenir institucionalmente?
  \end{itemize}
\end{itemize}

\subsection{Recapitulación}\label{recapitulaciuxf3n}

\begin{itemize}
\tightlist
\item
  \textbf{Si k=2}: Las titulaciones se agrupan principalmente en dos
  categorías amplias. Los 4 cuadrantes refinen estos dos clusters en
  distinciones más sutiles.
\item
  \textbf{Si k≥3}: Existen subgrupos adicionales dentro de las
  categorías principales.
\item
  \textbf{En ambos casos}: Los 4 cuadrantes enriquecen el análisis
  proporcionando un marco estratégico independiente del clustering
  automático.
\end{itemize}

\textbf{Conclusión}: No es una contradicción, sino \textbf{dos lentes
complementarios} para entender el mismo fenómeno: uno
estadístico-exploratorio (clusters) y otro operativo-estratégico
(cuadrantes).

\begin{center}\rule{0.5\linewidth}{0.5pt}\end{center}

\section{7. Interpretación de
Clústeres}\label{interpretaciuxf3n-de-cluxfasteres}

\begin{Shaded}
\begin{Highlighting}[]
\CommentTok{\# Ordenar titulaciones por clúster para mejor visualización}
\ControlFlowTok{for}\NormalTok{ (c }\ControlFlowTok{in} \DecValTok{1}\SpecialCharTok{:}\NormalTok{k\_optimal) \{}
  \FunctionTok{cat}\NormalTok{(}\StringTok{"}\SpecialCharTok{\textbackslash{}n}\StringTok{"}\NormalTok{, }\FunctionTok{strrep}\NormalTok{(}\StringTok{"="}\NormalTok{, }\DecValTok{80}\NormalTok{), }\StringTok{"}\SpecialCharTok{\textbackslash{}n}\StringTok{"}\NormalTok{)}
  \FunctionTok{cat}\NormalTok{(}\StringTok{"CLÚSTER"}\NormalTok{, c, }\StringTok{"}\SpecialCharTok{\textbackslash{}n}\StringTok{"}\NormalTok{)}
  \FunctionTok{cat}\NormalTok{(}\FunctionTok{strrep}\NormalTok{(}\StringTok{"="}\NormalTok{, }\DecValTok{80}\NormalTok{), }\StringTok{"}\SpecialCharTok{\textbackslash{}n}\StringTok{"}\NormalTok{)}
  
\NormalTok{  cluster\_data }\OtherTok{\textless{}{-}}\NormalTok{ titulaciones\_clustered }\SpecialCharTok{\%\textgreater{}\%}
    \FunctionTok{filter}\NormalTok{(Cluster }\SpecialCharTok{==}\NormalTok{ c) }\SpecialCharTok{\%\textgreater{}\%}
    \FunctionTok{arrange}\NormalTok{(}\FunctionTok{desc}\NormalTok{(Tasa\_Rendimiento))}
  
  \FunctionTok{cat}\NormalTok{(}\StringTok{"Número de titulaciones:"}\NormalTok{, }\FunctionTok{nrow}\NormalTok{(cluster\_data), }\StringTok{"}\SpecialCharTok{\textbackslash{}n}\StringTok{"}\NormalTok{)}
  \FunctionTok{cat}\NormalTok{(}\StringTok{"Rendimiento promedio:"}\NormalTok{, }
      \FunctionTok{round}\NormalTok{(}\FunctionTok{mean}\NormalTok{(cluster\_data}\SpecialCharTok{$}\NormalTok{Tasa\_Rendimiento), }\DecValTok{4}\NormalTok{), }\StringTok{"}\SpecialCharTok{\textbackslash{}n}\StringTok{"}\NormalTok{)}
  \FunctionTok{cat}\NormalTok{(}\StringTok{"Abandono promedio:"}\NormalTok{, }
      \FunctionTok{round}\NormalTok{(}\FunctionTok{mean}\NormalTok{(cluster\_data}\SpecialCharTok{$}\NormalTok{Tasa\_Abandono), }\DecValTok{4}\NormalTok{), }\StringTok{"}\SpecialCharTok{\textbackslash{}n\textbackslash{}n}\StringTok{"}\NormalTok{)}
  
  \FunctionTok{cat}\NormalTok{(}\StringTok{"Titulaciones incluidas:}\SpecialCharTok{\textbackslash{}n}\StringTok{"}\NormalTok{)}
  \ControlFlowTok{for}\NormalTok{ (i }\ControlFlowTok{in} \DecValTok{1}\SpecialCharTok{:}\FunctionTok{nrow}\NormalTok{(cluster\_data)) \{}
    \FunctionTok{cat}\NormalTok{(}\FunctionTok{sprintf}\NormalTok{(}
      \StringTok{"\%2d. \%s}\SpecialCharTok{\textbackslash{}n}\StringTok{   {-} Rendimiento: \%.1f\%\% | Abandono: \%.1f\%\%}\SpecialCharTok{\textbackslash{}n}\StringTok{"}\NormalTok{,}
\NormalTok{      i,}
\NormalTok{      cluster\_data}\SpecialCharTok{$}\NormalTok{Titulacion[i],}
\NormalTok{      cluster\_data}\SpecialCharTok{$}\NormalTok{Tasa\_Rendimiento[i] }\SpecialCharTok{*} \DecValTok{100}\NormalTok{,}
\NormalTok{      cluster\_data}\SpecialCharTok{$}\NormalTok{Tasa\_Abandono[i] }\SpecialCharTok{*} \DecValTok{100}
\NormalTok{    ))}
\NormalTok{  \}}
\NormalTok{\}}
\end{Highlighting}
\end{Shaded}

\begin{verbatim}
## 
##  ================================================================================ 
## CLÚSTER 1 
## ================================================================================ 
## Número de titulaciones: 59 
## Rendimiento promedio: 0.5927 
## Abandono promedio: 0.3106 
## 
## Titulaciones incluidas:
##  1. Máster Un. en Asesoría Jurídico-Mercantil, Fiscal y Laboral
##    - Rendimiento: 79.3% | Abandono: 30.8%
##  2. Grado en Marketing e Investigación de Mercados
##    - Rendimiento: 75.4% | Abandono: 26.5%
##  3. Máster Universitario en Derecho Constitucional
##    - Rendimiento: 73.0% | Abandono: 27.9%
##  4. Máster Universitario en Estudios Hispánicos Superiores
##    - Rendimiento: 70.8% | Abandono: 28.9%
##  5. Grado en Ingeniería Informática-Ingeniería del Software
##    - Rendimiento: 70.7% | Abandono: 26.3%
##  6. Grado en Estudios Ingleses
##    - Rendimiento: 70.2% | Abandono: 25.3%
##  7. Grado en Ingeniería de la Salud por la U. Málaga y la U.Sevilla
##    - Rendimiento: 68.6% | Abandono: 24.4%
##  8. Grado en Filosofía
##    - Rendimiento: 67.3% | Abandono: 44.3%
##  9. Grado en Filología Clásica
##    - Rendimiento: 67.2% | Abandono: 55.9%
## 10. Grado en Matemáticas
##    - Rendimiento: 66.7% | Abandono: 23.2%
## 11. Grado en Lengua y Literatura Alemanas
##    - Rendimiento: 66.6% | Abandono: 37.0%
## 12. Grado en Geografia y Gestión del Territorio
##    - Rendimiento: 66.4% | Abandono: 33.2%
## 13. Máster Universitario en Peritación y Reparación de Edificios
##    - Rendimiento: 66.4% | Abandono: 26.8%
## 14. Grado en Historia
##    - Rendimiento: 66.2% | Abandono: 26.5%
## 15. Grado en Geografía y Gestión del Territorio
##    - Rendimiento: 66.0% | Abandono: 42.7%
## 16. Grado en Ingeniería de Organización Industrial por la U. M. y U.S.
##    - Rendimiento: 65.9% | Abandono: 24.1%
## 17. Grado en Estudios Árabes e Islámicos
##    - Rendimiento: 65.4% | Abandono: 43.0%
## 18. Grado en Relaciones Laborales y Recursos Humanos
##    - Rendimiento: 65.1% | Abandono: 22.3%
## 19. Grado en Física
##    - Rendimiento: 64.6% | Abandono: 26.8%
## 20. Grado en Ingeniería de Materiales
##    - Rendimiento: 64.5% | Abandono: 23.5%
## 21. Grado en Turismo
##    - Rendimiento: 64.5% | Abandono: 29.9%
## 22. Grado en Filología Hispánica
##    - Rendimiento: 62.6% | Abandono: 32.3%
## 23. Grado en Antropología Social y Cultural
##    - Rendimiento: 62.5% | Abandono: 45.2%
## 24. M.U. en Lógica, Computación e Inteligencia Artificial
##    - Rendimiento: 62.5% | Abandono: 18.0%
## 25. Grado en Ingeniería de Tecnologías Industriales
##    - Rendimiento: 62.4% | Abandono: 16.8%
## 26. Grado en Historia del Arte
##    - Rendimiento: 62.1% | Abandono: 39.2%
## 27. Grado en Ingeniería de Organización Industrial por la U. Málaga y U.Sevilla
##    - Rendimiento: 62.0% | Abandono: 27.0%
## 28. Grado en Estudios Franceses
##    - Rendimiento: 62.0% | Abandono: 39.5%
## 29. Grado en Química
##    - Rendimiento: 62.0% | Abandono: 27.0%
## 30. Grado en Ingeniería en Diseño Industrial y Desarrollo del Producto
##    - Rendimiento: 61.1% | Abandono: 22.6%
## 31. Grado en Economía
##    - Rendimiento: 60.5% | Abandono: 32.0%
## 32. Grado en Ingeniería de la Energía por la US y UMA
##    - Rendimiento: 58.3% | Abandono: 37.3%
## 33. Grado en Administración y Dirección de Empresas
##    - Rendimiento: 57.9% | Abandono: 37.1%
## 34. Máster Unv. en Microelectrónica: Diseño y Aplicaciones de Sistemas
##    - Rendimiento: 57.7% | Abandono: 35.8%
## 35. Máster Universitario en Sistemas de Energía Térmica
##    - Rendimiento: 56.7% | Abandono: 31.7%
## 36. Máster Universitario en Economía y Desarrollo
##    - Rendimiento: 56.5% | Abandono: 25.9%
## 37. M.U. en Diseño Avanzado en Ingeniería Mecánica
##    - Rendimiento: 56.4% | Abandono: 32.2%
## 38. Máster Universitario en Organización Industrial y Gestión de Empresas
##    - Rendimiento: 56.4% | Abandono: 10.4%
## 39. Grado en Ingeniería Informática-Tecnologías Informáticas
##    - Rendimiento: 56.4% | Abandono: 45.3%
## 40. Grado en Ingeniería de la Energía por la U.S y la U. Málaga
##    - Rendimiento: 56.0% | Abandono: 23.6%
## 41. Grado en Ingeniería de la Energía por la Un. de Sevilla y la Un. de Málaga
##    - Rendimiento: 55.1% | Abandono: 28.1%
## 42. M.U. en Sistemas de Energía Eléctrica
##    - Rendimiento: 54.0% | Abandono: 21.7%
## 43. Grado en Ingeniería de las Tecnologías de Telecomunicación
##    - Rendimiento: 53.1% | Abandono: 32.7%
## 44. Máster Universitario en Sistemas de Energía Eléctrica
##    - Rendimiento: 52.8% | Abandono: 14.8%
## 45. M.U. en Microelectrónica: Diseño y Aplicaciones de Sistemas
##    - Rendimiento: 52.5% | Abandono: 31.9%
## 46. Grado en Ingeniería Química
##    - Rendimiento: 52.4% | Abandono: 36.3%
## 47. Grado en Ingeniería Mecánica
##    - Rendimiento: 51.7% | Abandono: 30.9%
## 48. Máster Universitario en Arquitectura
##    - Rendimiento: 51.3% | Abandono: 6.4%
## 49. Grado en Ingeniería Civil
##    - Rendimiento: 51.2% | Abandono: 45.0%
## 50. Máster Universitario en Comunicación y Cultura
##    - Rendimiento: 51.1% | Abandono: 16.3%
## 51. Grado en Ingeniería Informática-Ingeniería de Computadores
##    - Rendimiento: 49.7% | Abandono: 59.2%
## 52. Grado en Finanzas y Contabilidad
##    - Rendimiento: 48.9% | Abandono: 49.3%
## 53. Grado en Ingeniería Eléctrica
##    - Rendimiento: 46.3% | Abandono: 44.0%
## 54. Grado en Ingeniería Química Industrial
##    - Rendimiento: 45.6% | Abandono: 44.5%
## 55. Máster Universitario en Comunicación Institucional y Política
##    - Rendimiento: 45.1% | Abandono: 14.6%
## 56. Máster Universitario en Especialización Profesional en Farmacia
##    - Rendimiento: 45.0% | Abandono: 17.0%
## 57. Máster Universitario en Diseño Avanzado en Ingeniería Mecánica
##    - Rendimiento: 44.8% | Abandono: 28.2%
## 58. Grado en Ingeniería Electrónica Industrial
##    - Rendimiento: 43.5% | Abandono: 46.1%
## 59. M.U. en Microelectrónica: Dño. y Apps.de Sistemas M./Nano.
##    - Rendimiento: 39.9% | Abandono: 37.9%
## 
##  ================================================================================ 
## CLÚSTER 2 
## ================================================================================ 
## Número de titulaciones: 139 
## Rendimiento promedio: 0.8108 
## Abandono promedio: 0.1071 
## 
## Titulaciones incluidas:
##  1. M.U. en Estudios Hispánicos Superiores (R.D.1393/07)
##    - Rendimiento: 113.3% | Abandono: 7.2%
##  2. M.U. en Filosofía y Cultura Moderna (R.D.1393/2007)
##    - Rendimiento: 103.9% | Abandono: 10.3%
##  3. M.U. Profesorado de E.S.O. y Bachillerato, F. P. y E.de Idiomas
##    - Rendimiento: 98.6% | Abandono: 2.9%
##  4. Grado en Educación Infantil
##    - Rendimiento: 97.4% | Abandono: 5.4%
##  5. M.U. en Profesorado de E.S.O y Bachillerato, FP y E.Idiomas
##    - Rendimiento: 95.6% | Abandono: 3.3%
##  6. Máster Univers. en Profesorado de E.S.O y Bachillerato, FP y E.Idiomas
##    - Rendimiento: 94.0% | Abandono: 3.0%
##  7. M.U. en Estudios Avanzados en Química (R.D.1393/07)
##    - Rendimiento: 94.0% | Abandono: 0.0%
##  8. M.U. en Ciencia y Tecnología de Nuevos Materiales (R.D.1393/07)
##    - Rendimiento: 93.9% | Abandono: 10.5%
##  9. Grado en Criminología
##    - Rendimiento: 93.7% | Abandono: 11.4%
## 10. M.U. en Psicopedagogía
##    - Rendimiento: 92.8% | Abandono: 2.7%
## 11. M.U. en En.del Español como L. Extranjera y de Otras L. Modernas
##    - Rendimiento: 91.8% | Abandono: 9.5%
## 12. Grado en Biomedicina Básica y Experimental
##    - Rendimiento: 91.8% | Abandono: 17.1%
## 13. Grado en Pedagogía
##    - Rendimiento: 91.5% | Abandono: 12.6%
## 14. Máster Univ. en Psicología de la Intervención Social y Comunitaria
##    - Rendimiento: 89.9% | Abandono: 4.5%
## 15. Grado en Publicidad y Relaciones Públicas
##    - Rendimiento: 89.7% | Abandono: 12.0%
## 16. Grado en Educación Primaria
##    - Rendimiento: 89.5% | Abandono: 10.2%
## 17. M.U. en Estudios Avanzados en Química
##    - Rendimiento: 88.9% | Abandono: 3.3%
## 18. M.U. en Dirección, Evaluación y Calidad de las Instituc. de For
##    - Rendimiento: 88.8% | Abandono: 3.7%
## 19. Grado en Comunicación Audiovisual
##    - Rendimiento: 88.8% | Abandono: 11.6%
## 20. Grado en Óptica y Optometría
##    - Rendimiento: 88.7% | Abandono: 12.5%
## 21. M.U. en Formación y Orientación para el Trabajo
##    - Rendimiento: 88.6% | Abandono: 1.4%
## 22. M.U. en Biología Avanzada: Investigación y Aplicación
##    - Rendimiento: 88.6% | Abandono: 8.4%
## 23. Máster Universitario en Psicopedagogía
##    - Rendimiento: 88.5% | Abandono: 3.5%
## 24. M.U. en Estudios Americanos (R.D.1393/2007)
##    - Rendimiento: 88.4% | Abandono: 9.8%
## 25. M.U. en Ciencia y Tecnología de Nuevos Materiales
##    - Rendimiento: 88.2% | Abandono: 12.7%
## 26. M.U. en Seguridad Integral en Edificación (R.D.1393/07)
##    - Rendimiento: 88.2% | Abandono: 0.0%
## 27. Grado en Ciencias de la Actividad Física y del Deporte
##    - Rendimiento: 88.1% | Abandono: 10.5%
## 28. M.U. en Seguridad Integral en Edificación
##    - Rendimiento: 88.1% | Abandono: 12.7%
## 29. Máster Universitario en Biología Avanzada: Investigación y Aplicación
##    - Rendimiento: 87.9% | Abandono: 8.1%
## 30. Máster Universitario en Formación y Orientación para el Trabajo
##    - Rendimiento: 87.3% | Abandono: 3.6%
## 31. Grado en Bioquímica por la U. de Sevilla y U. de Málaga
##    - Rendimiento: 86.9% | Abandono: 13.6%
## 32. Grado en Bellas Artes
##    - Rendimiento: 86.9% | Abandono: 18.2%
## 33. Grado en Psicología
##    - Rendimiento: 86.9% | Abandono: 12.9%
## 34. M.U. en Dirección, Evaluación y Calidad de las Instituciones de For.
##    - Rendimiento: 86.8% | Abandono: 15.2%
## 35. M.U. en Estudios Americanos
##    - Rendimiento: 86.7% | Abandono: 8.2%
## 36. Grado en Fisioterapia
##    - Rendimiento: 86.7% | Abandono: 26.1%
## 37. Grado en Bioquímica por la U.S. y UMA
##    - Rendimiento: 86.4% | Abandono: 10.4%
## 38. Grado en Enfermería
##    - Rendimiento: 86.4% | Abandono: 21.5%
## 39. M.U.en Enseñanza del Español como Lengua Extranjera y Otras Leng.Mod.
##    - Rendimiento: 86.1% | Abandono: 5.6%
## 40. Grado en Estudios de Asia Oriental por la US y UMA
##    - Rendimiento: 86.1% | Abandono: 17.3%
## 41. M.U. en Psicología de la Intervención Social y Comunitaria
##    - Rendimiento: 86.1% | Abandono: 3.4%
## 42. Máster Universitario en Ciencia y Tecnología de Nuevos Materiales
##    - Rendimiento: 85.8% | Abandono: 10.4%
## 43. M.U. en Estudios Lingüísticos, Literarios y Culturales por la U.S.
##    - Rendimiento: 85.8% | Abandono: 7.0%
## 44. M. U. en Estudios Históricos Avanzados
##    - Rendimiento: 85.7% | Abandono: 12.1%
## 45. M.U. en Estudios de Género y Desarrollo Profesional
##    - Rendimiento: 85.7% | Abandono: 11.8%
## 46. M.U. en Artes del Espectáculo Vivo (R.D.1393/07)
##    - Rendimiento: 85.4% | Abandono: 5.6%
## 47. M.U. en Comunicación y Cultura (R.D.1393/07)
##    - Rendimiento: 85.0% | Abandono: 13.8%
## 48. Máster Unv. en Dirección, Evaluación y Calidad de las Instituc. de For
##    - Rendimiento: 84.5% | Abandono: 6.9%
## 49. M.U.en Seguridad Integral en la Industria y Prevención Riesgos Laborales
##    - Rendimiento: 84.5% | Abandono: 9.2%
## 50. Máster Unv.en Patrimonio Artístico Andaluz y su Proyección Iberoameric.
##    - Rendimiento: 84.4% | Abandono: 9.7%
## 51. Máster Universitario en Estudios Avanzados en Química
##    - Rendimiento: 84.3% | Abandono: 5.6%
## 52. Grado en Conservación y Restauración de Bienes Culturales
##    - Rendimiento: 84.3% | Abandono: 25.9%
## 53. M.U. en Comunicación Institucional y Política
##    - Rendimiento: 84.2% | Abandono: 11.2%
## 54. M.U. en Estudios Históricos Avanzados
##    - Rendimiento: 84.2% | Abandono: 15.5%
## 55. Grado en Estudios de Asia Oriental por la U.S. y la U.de Málaga
##    - Rendimiento: 84.0% | Abandono: 21.2%
## 56. Máster Universitario en Filosofía y Cultura Moderna
##    - Rendimiento: 84.0% | Abandono: 9.8%
## 57. Grado en Bioquímica por la Universidad de Sevilla y Universidad de Málaga
##    - Rendimiento: 83.9% | Abandono: 11.3%
## 58. M.U. en Migraciones Internacionales, Salud y Bienestar: Modelos y Estrategias de Intervención
##    - Rendimiento: 83.7% | Abandono: 11.0%
## 59. M.U.en Migraciones Internacionales, Salud y Bienestar: Mod. y Est. de Inter
##    - Rendimiento: 83.3% | Abandono: 2.8%
## 60. M.U. en Patrimonio Artístico Andaluz y su Proy. Iberoamericana
##    - Rendimiento: 83.2% | Abandono: 6.7%
## 61. M.U. en Consultoría Laboral
##    - Rendimiento: 82.9% | Abandono: 5.2%
## 62. Grado en Estudios de Asia Oriental por la Unv.de Sevilla y la Unv.de Málaga
##    - Rendimiento: 82.9% | Abandono: 17.8%
## 63. M.U. en Documentos y Libros. Archivos y Bibliotecas
##    - Rendimiento: 82.8% | Abandono: 1.5%
## 64. Máster Universitario en Estudios Históricos Avanzados
##    - Rendimiento: 82.8% | Abandono: 14.0%
## 65. M.U. en Especialización Profesional en Farmacia
##    - Rendimiento: 82.8% | Abandono: 7.5%
## 66. M.U. en Relaciones Jurídico-Privadas (R.D.1393/07)
##    - Rendimiento: 82.7% | Abandono: 3.6%
## 67. M. U. en Gestión del Territorio, Instrumentos y Técnicas de Intervenc.
##    - Rendimiento: 82.4% | Abandono: 5.6%
## 68. M.U. en Derecho Constitucional (R.D.1393/07)
##    - Rendimiento: 82.1% | Abandono: 19.2%
## 69. Máster Unv. en Investigación Médica: Clínica y Experimental
##    - Rendimiento: 81.6% | Abandono: 8.8%
## 70. M.U. en Gestión Integral de la Edificación (R.D.1393/07)
##    - Rendimiento: 81.4% | Abandono: 7.1%
## 71. Máster Universitario en Estudios de Género y Desarrollo Profesional
##    - Rendimiento: 81.3% | Abandono: 10.0%
## 72. M.U. en Ciudad y Arquitectura Sostenibles
##    - Rendimiento: 81.0% | Abandono: 8.8%
## 73. M.U. en Investigación Médica: Clínica y Experimental
##    - Rendimiento: 81.0% | Abandono: 14.1%
## 74. Máster Universitario en Estudios Americanos
##    - Rendimiento: 81.0% | Abandono: 6.4%
## 75. M.U. en Seguridad Integral en la Industria y Prev. RR.LL.
##    - Rendimiento: 81.0% | Abandono: 7.5%
## 76. M.U. en Gestión Integral de la Edificación
##    - Rendimiento: 80.9% | Abandono: 8.3%
## 77. M.U. en Estudios de Género y Desarrollo Profesional (RD.07)
##    - Rendimiento: 80.8% | Abandono: 9.2%
## 78. M.U. en Asesoría Jurídico-Mercantil, Fiscal y Laboral
##    - Rendimiento: 80.7% | Abandono: 5.5%
## 79. Máster Universitario en Consultoría Laboral
##    - Rendimiento: 80.5% | Abandono: 14.4%
## 80. M.U. en Dirección y Planificación del Turismo
##    - Rendimiento: 80.0% | Abandono: 11.0%
## 81. M.U. en Patrimonio Artístico Andaluz y su Proyección Iberoameric.
##    - Rendimiento: 80.0% | Abandono: 10.8%
## 82. M.U. en Estudios Lingüísticos, Literarios y Culturales
##    - Rendimiento: 79.6% | Abandono: 8.4%
## 83. M.U. en Arte: Idea y Producción
##    - Rendimiento: 79.5% | Abandono: 7.7%
## 84. Grado en Periodismo
##    - Rendimiento: 79.3% | Abandono: 14.8%
## 85. Grado en Podología
##    - Rendimiento: 79.3% | Abandono: 28.2%
## 86. M.U. en Dirección y Planificación del Turismo (R.D.1393/07)
##    - Rendimiento: 79.2% | Abandono: 5.4%
## 87. Máster Universitario en Documentos y Libros. Archivos y Bibliotecas
##    - Rendimiento: 79.2% | Abandono: 4.1%
## 88. M.U. en Comunicación y Cultura
##    - Rendimiento: 78.9% | Abandono: 6.0%
## 89. M.U. en Gestión del Territorio. Instrumentos y Técnicas de Intervenc
##    - Rendimiento: 78.9% | Abandono: 10.6%
## 90. M.U. en Comunicación Institucional y Política (R.D.1393/2007)
##    - Rendimiento: 78.8% | Abandono: 6.0%
## 91. M.U. en Nuevas Tendencias Asistenciales en CC. de la Salud
##    - Rendimiento: 78.7% | Abandono: 5.3%
## 92. M.U. en Escritura Creativa
##    - Rendimiento: 78.6% | Abandono: 12.9%
## 93. M.U. en Derecho Constitucional
##    - Rendimiento: 78.3% | Abandono: 10.6%
## 94. Máster Universitario en Relaciones Jurídico-Privadas
##    - Rendimiento: 77.9% | Abandono: 23.7%
## 95. M.U. en Nuevas Tendencias Asistenciales en Ciencias de la Salud
##    - Rendimiento: 77.6% | Abandono: 8.4%
## 96. Máster Unv.en Patrimonio Artístico Andaluz y su Proyección Iberoamericana
##    - Rendimiento: 77.5% | Abandono: 5.1%
## 97. Máster Universitario en Estudios Lingüísticos, Literarios y Culturales
##    - Rendimiento: 77.4% | Abandono: 8.3%
## 98. M.U. en Escritura Creativa (R.D.1393/07)
##    - Rendimiento: 76.6% | Abandono: 14.1%
## 99. M.U. en Arte: Idea y Producción (R.D.1393/07)
##    - Rendimiento: 76.3% | Abandono: 7.7%
## 100. Máster Universitario en Ciudad y Arquitectura Sostenibles
##    - Rendimiento: 76.3% | Abandono: 6.8%
## 101. M.U. en Gestión Estratégica y Negocios Internacionales
##    - Rendimiento: 76.3% | Abandono: 10.9%
## 102. Grado en Biología
##    - Rendimiento: 76.1% | Abandono: 15.7%
## 103. M.U. en Artes del Espectáculo Vivo
##    - Rendimiento: 75.7% | Abandono: 17.1%
## 104. Máster Universitario en Seguridad Integral en Edificación
##    - Rendimiento: 75.6% | Abandono: 25.3%
## 105. M.U. en Estudios Hispánicos Superiores
##    - Rendimiento: 75.5% | Abandono: 13.5%
## 106. Máster Universitario en Gestión Estratégica y Negocios Internacionales
##    - Rendimiento: 75.2% | Abandono: 15.1%
## 107. Máster Universitario en Escritura Creativa
##    - Rendimiento: 73.9% | Abandono: 5.6%
## 108. M.U. en Arqueología por la US y UGR
##    - Rendimiento: 73.7% | Abandono: 10.5%
## 109. Máster Universitario en Lógica, Computación e Inteligencia Artificial
##    - Rendimiento: 73.5% | Abandono: 12.5%
## 110. Máster Universitario en Gestión Integral de la Edificación
##    - Rendimiento: 73.4% | Abandono: 16.0%
## 111. M.U. en Antropología: Gestión de la Diver. Cul., el Patr. y el Des.
##    - Rendimiento: 73.4% | Abandono: 17.6%
## 112. Máster Universitario en Arte: Idea y Producción
##    - Rendimiento: 73.2% | Abandono: 4.7%
## 113. Grado en Ingeniería de la Salud por la Univ. de Málaga y la Univ.de Sevilla
##    - Rendimiento: 73.2% | Abandono: 14.8%
## 114. M.U. en Arquitectura
##    - Rendimiento: 72.4% | Abandono: 4.2%
## 115. M.U. en Matemáticas
##    - Rendimiento: 72.3% | Abandono: 11.9%
## 116. Grado en Derecho
##    - Rendimiento: 72.1% | Abandono: 14.6%
## 117. Máster Universitario en Dirección y Planificación del Turismo
##    - Rendimiento: 72.0% | Abandono: 20.5%
## 118. Máster Universitario en Artes del Espectáculo Vivo
##    - Rendimiento: 71.6% | Abandono: 12.2%
## 119. M.U. en Arqueología (R.D.1393/07)
##    - Rendimiento: 71.4% | Abandono: 16.3%
## 120. Máster Universitario en Física Nuclear por la UAM,UCM,UB,UGR,USAL y la US
##    - Rendimiento: 71.3% | Abandono: 9.9%
## 121. M.U. en Física Nuclear
##    - Rendimiento: 71.1% | Abandono: 8.9%
## 122. M.U. en Filosofía y Cultura Moderna
##    - Rendimiento: 70.9% | Abandono: 11.1%
## 123. Máster Universitario en Matemáticas
##    - Rendimiento: 70.9% | Abandono: 10.6%
## 124. M.U. en Relaciones Jurídicas-Privadas
##    - Rendimiento: 69.9% | Abandono: 18.1%
## 125. M.U. en Economía y Desarrollo (R.D.1393/07)
##    - Rendimiento: 69.3% | Abandono: 17.3%
## 126. Máster Unv. en Gestión del Territorio. Instrumentos y Técnicas de Intervenc
##    - Rendimiento: 69.1% | Abandono: 18.3%
## 127. M.U. en Economía y Desarrollo
##    - Rendimiento: 69.0% | Abandono: 19.6%
## 128. M.U. en Física Nuclear (Máster Interuniversitario) (R.D.1393/07)
##    - Rendimiento: 68.9% | Abandono: 6.2%
## 129. Grado en Ingeniería de la Salud por la UMA y US
##    - Rendimiento: 68.9% | Abandono: 17.4%
## 130. M.U. en Peritación y Reparación de Edificios
##    - Rendimiento: 68.6% | Abandono: 4.3%
## 131. Grado en Ingeniería Electrónica, Robótica y Mecatrónica (UMA-US)
##    - Rendimiento: 67.8% | Abandono: 17.6%
## 132. Grado en Ingeniería Aeroespacial
##    - Rendimiento: 66.3% | Abandono: 16.0%
## 133. Máster Universitario en Asesoría Jurídico-Mercantil, Fiscal y Laboral
##    - Rendimiento: 66.2% | Abandono: 6.2%
## 134. M.U. en Antropología: Gestión de la Diversidad Cultural, el Pat.y el Des.
##    - Rendimiento: 66.0% | Abandono: 16.6%
## 135. M.U. en Sistemas de Energía Térmica
##    - Rendimiento: 65.6% | Abandono: 10.6%
## 136. Grado en Ingeniería de Organización Industrial por la UMA y US
##    - Rendimiento: 65.2% | Abandono: 14.5%
## 137. M.U. en Organización Industrial y Gestión de Empresas
##    - Rendimiento: 64.7% | Abandono: 12.9%
## 138. Máster Universitario en Arqueología por la U.de Sevilla y U.de Granada
##    - Rendimiento: 63.2% | Abandono: 10.6%
## 139. Máster Univ.en Nuevas Tendencias Asistenciales en CC. de la Salud
##    - Rendimiento: 59.8% | Abandono: 7.6%
\end{verbatim}

\begin{center}\rule{0.5\linewidth}{0.5pt}\end{center}

\section{8. Análisis de Calidad del
Clustering}\label{anuxe1lisis-de-calidad-del-clustering}

\begin{Shaded}
\begin{Highlighting}[]
\CommentTok{\# Calcular métricas de calidad}
\NormalTok{silhouette\_coef }\OtherTok{\textless{}{-}} \FunctionTok{silhouette}\NormalTok{(kmeans\_result}\SpecialCharTok{$}\NormalTok{cluster, }\FunctionTok{dist}\NormalTok{(features\_normalized))}
\FunctionTok{cat}\NormalTok{(}\StringTok{"Coeficiente de Silhueta Global:"}\NormalTok{, }
    \FunctionTok{round}\NormalTok{(}\FunctionTok{mean}\NormalTok{(silhouette\_coef[, }\StringTok{"sil\_width"}\NormalTok{]), }\DecValTok{4}\NormalTok{), }\StringTok{"}\SpecialCharTok{\textbackslash{}n\textbackslash{}n}\StringTok{"}\NormalTok{)}
\end{Highlighting}
\end{Shaded}

\begin{verbatim}
## Coeficiente de Silhueta Global: 0.5382
\end{verbatim}

\begin{Shaded}
\begin{Highlighting}[]
\CommentTok{\# Davies{-}Bouldin Index (menor es mejor)}
\NormalTok{db\_index }\OtherTok{\textless{}{-}} \DecValTok{0}
\ControlFlowTok{for}\NormalTok{ (i }\ControlFlowTok{in} \DecValTok{1}\SpecialCharTok{:}\NormalTok{k\_optimal) \{}
\NormalTok{  cluster\_i }\OtherTok{\textless{}{-}} \FunctionTok{which}\NormalTok{(kmeans\_result}\SpecialCharTok{$}\NormalTok{cluster }\SpecialCharTok{==}\NormalTok{ i)}
\NormalTok{  center\_i }\OtherTok{\textless{}{-}}\NormalTok{ kmeans\_result}\SpecialCharTok{$}\NormalTok{centers[i, ]}
  
\NormalTok{  diameter\_i }\OtherTok{\textless{}{-}} \FunctionTok{max}\NormalTok{(}\FunctionTok{sapply}\NormalTok{(cluster\_i, }\ControlFlowTok{function}\NormalTok{(x) \{}
    \FunctionTok{dist}\NormalTok{(}\FunctionTok{rbind}\NormalTok{(features\_normalized[x, ], center\_i))}
\NormalTok{  \}))}
  
  \ControlFlowTok{for}\NormalTok{ (j }\ControlFlowTok{in} \DecValTok{1}\SpecialCharTok{:}\NormalTok{k\_optimal) \{}
    \ControlFlowTok{if}\NormalTok{ (i }\SpecialCharTok{!=}\NormalTok{ j) \{}
\NormalTok{      center\_j }\OtherTok{\textless{}{-}}\NormalTok{ kmeans\_result}\SpecialCharTok{$}\NormalTok{centers[j, ]}
\NormalTok{      distance\_ij }\OtherTok{\textless{}{-}} \FunctionTok{dist}\NormalTok{(}\FunctionTok{rbind}\NormalTok{(center\_i, center\_j))}
      
\NormalTok{      ratio }\OtherTok{\textless{}{-}}\NormalTok{ (diameter\_i }\SpecialCharTok{+} \FunctionTok{max}\NormalTok{(}\FunctionTok{sapply}\NormalTok{(}\FunctionTok{which}\NormalTok{(kmeans\_result}\SpecialCharTok{$}\NormalTok{cluster }\SpecialCharTok{==}\NormalTok{ j), }\ControlFlowTok{function}\NormalTok{(x) \{}
        \FunctionTok{dist}\NormalTok{(}\FunctionTok{rbind}\NormalTok{(features\_normalized[x, ], center\_j))}
\NormalTok{      \}))) }\SpecialCharTok{/}\NormalTok{ distance\_ij}
      
\NormalTok{      db\_index }\OtherTok{\textless{}{-}} \FunctionTok{max}\NormalTok{(db\_index, ratio)}
\NormalTok{    \}}
\NormalTok{  \}}
\NormalTok{\}}

\FunctionTok{cat}\NormalTok{(}\StringTok{"Davies{-}Bouldin Index:"}\NormalTok{, }\FunctionTok{round}\NormalTok{(db\_index, }\DecValTok{4}\NormalTok{), }\StringTok{"}\SpecialCharTok{\textbackslash{}n}\StringTok{"}\NormalTok{)}
\end{Highlighting}
\end{Shaded}

\begin{verbatim}
## Davies-Bouldin Index: 2.0744
\end{verbatim}

\begin{Shaded}
\begin{Highlighting}[]
\FunctionTok{cat}\NormalTok{(}\StringTok{"(Valores menores indican mejor separación entre clústeres)}\SpecialCharTok{\textbackslash{}n\textbackslash{}n}\StringTok{"}\NormalTok{)}
\end{Highlighting}
\end{Shaded}

\begin{verbatim}
## (Valores menores indican mejor separación entre clústeres)
\end{verbatim}

\begin{Shaded}
\begin{Highlighting}[]
\CommentTok{\# Varianza explicada}
\FunctionTok{cat}\NormalTok{(}\StringTok{"Varianza intra{-}cluster (WSS):"}\NormalTok{, }\FunctionTok{round}\NormalTok{(kmeans\_result}\SpecialCharTok{$}\NormalTok{tot.withinss, }\DecValTok{4}\NormalTok{), }\StringTok{"}\SpecialCharTok{\textbackslash{}n}\StringTok{"}\NormalTok{)}
\end{Highlighting}
\end{Shaded}

\begin{verbatim}
## Varianza intra-cluster (WSS): 160.8865
\end{verbatim}

\begin{Shaded}
\begin{Highlighting}[]
\FunctionTok{cat}\NormalTok{(}\StringTok{"Varianza total:"}\NormalTok{, }\FunctionTok{round}\NormalTok{(}\FunctionTok{sum}\NormalTok{(}\FunctionTok{apply}\NormalTok{(features\_normalized, }\DecValTok{2}\NormalTok{, var)) }\SpecialCharTok{*} \FunctionTok{nrow}\NormalTok{(features\_normalized), }\DecValTok{4}\NormalTok{), }\StringTok{"}\SpecialCharTok{\textbackslash{}n}\StringTok{"}\NormalTok{)}
\end{Highlighting}
\end{Shaded}

\begin{verbatim}
## Varianza total: 396
\end{verbatim}

\begin{center}\rule{0.5\linewidth}{0.5pt}\end{center}

\section{9. Recomendaciones para Intervención
Institucional}\label{recomendaciones-para-intervenciuxf3n-institucional}

\begin{Shaded}
\begin{Highlighting}[]
\CommentTok{\# Calcular medianas para clasificación de cuadrantes}
\NormalTok{median\_rend }\OtherTok{\textless{}{-}} \FunctionTok{median}\NormalTok{(titulaciones\_clustered}\SpecialCharTok{$}\NormalTok{Tasa\_Rendimiento)}
\NormalTok{median\_aband }\OtherTok{\textless{}{-}} \FunctionTok{median}\NormalTok{(titulaciones\_clustered}\SpecialCharTok{$}\NormalTok{Tasa\_Abandono)}

\CommentTok{\# Crear variable de cuadrante para cada titulación}
\NormalTok{titulaciones\_quadrants }\OtherTok{\textless{}{-}}\NormalTok{ titulaciones\_clustered }\SpecialCharTok{\%\textgreater{}\%}
  \FunctionTok{mutate}\NormalTok{(}
    \AttributeTok{Cuadrante =} \FunctionTok{case\_when}\NormalTok{(}
\NormalTok{      Tasa\_Rendimiento }\SpecialCharTok{\textgreater{}=}\NormalTok{ median\_rend }\SpecialCharTok{\&}\NormalTok{ Tasa\_Abandono }\SpecialCharTok{\textless{}}\NormalTok{ median\_aband }\SpecialCharTok{\textasciitilde{}} 
        \StringTok{"Q1: EXITOSAS"}\NormalTok{,}
\NormalTok{      Tasa\_Rendimiento }\SpecialCharTok{\textgreater{}=}\NormalTok{ median\_rend }\SpecialCharTok{\&}\NormalTok{ Tasa\_Abandono }\SpecialCharTok{\textgreater{}=}\NormalTok{ median\_aband }\SpecialCharTok{\textasciitilde{}} 
        \StringTok{"Q2: ABANDONO SELECTIVO"}\NormalTok{,}
\NormalTok{      Tasa\_Rendimiento }\SpecialCharTok{\textless{}}\NormalTok{ median\_rend }\SpecialCharTok{\&}\NormalTok{ Tasa\_Abandono }\SpecialCharTok{\textless{}}\NormalTok{ median\_aband }\SpecialCharTok{\textasciitilde{}} 
        \StringTok{"Q3: PROBLEMÁTICO"}\NormalTok{,}
\NormalTok{      Tasa\_Rendimiento }\SpecialCharTok{\textless{}}\NormalTok{ median\_rend }\SpecialCharTok{\&}\NormalTok{ Tasa\_Abandono }\SpecialCharTok{\textgreater{}=}\NormalTok{ median\_aband }\SpecialCharTok{\textasciitilde{}} 
        \StringTok{"Q4: EN RIESGO"}
\NormalTok{    )}
\NormalTok{  )}

\FunctionTok{cat}\NormalTok{(}\StringTok{"}\SpecialCharTok{\textbackslash{}n}\StringTok{"}\NormalTok{)}
\end{Highlighting}
\end{Shaded}

\begin{Shaded}
\begin{Highlighting}[]
\FunctionTok{cat}\NormalTok{(}\StringTok{"CLASIFICACIÓN POR CUADRANTES (4 MATRICES) Y RECOMENDACIONES}\SpecialCharTok{\textbackslash{}n}\StringTok{"}\NormalTok{)}
\end{Highlighting}
\end{Shaded}

\begin{verbatim}
## CLASIFICACIÓN POR CUADRANTES (4 MATRICES) Y RECOMENDACIONES
\end{verbatim}

\begin{Shaded}
\begin{Highlighting}[]
\FunctionTok{cat}\NormalTok{(}\FunctionTok{strrep}\NormalTok{(}\StringTok{"="}\NormalTok{, }\DecValTok{130}\NormalTok{), }\StringTok{"}\SpecialCharTok{\textbackslash{}n\textbackslash{}n}\StringTok{"}\NormalTok{)}
\end{Highlighting}
\end{Shaded}

\begin{verbatim}
## ==================================================================================================================================
\end{verbatim}

\begin{Shaded}
\begin{Highlighting}[]
\CommentTok{\# Q1: EXITOSAS}
\NormalTok{q1\_data }\OtherTok{\textless{}{-}}\NormalTok{ titulaciones\_quadrants }\SpecialCharTok{\%\textgreater{}\%} \FunctionTok{filter}\NormalTok{(}\FunctionTok{str\_detect}\NormalTok{(Cuadrante, }\StringTok{"Q1"}\NormalTok{))}
\ControlFlowTok{if}\NormalTok{ (}\FunctionTok{nrow}\NormalTok{(q1\_data) }\SpecialCharTok{\textgreater{}} \DecValTok{0}\NormalTok{) \{}
  \FunctionTok{cat}\NormalTok{(}\StringTok{"Q1: TITULACIONES EXITOSAS | (Alto Rendimiento, Bajo Abandono)}\SpecialCharTok{\textbackslash{}n}\StringTok{"}\NormalTok{)}
  \FunctionTok{cat}\NormalTok{(}\FunctionTok{strrep}\NormalTok{(}\StringTok{"{-}"}\NormalTok{, }\DecValTok{130}\NormalTok{), }\StringTok{"}\SpecialCharTok{\textbackslash{}n}\StringTok{"}\NormalTok{)}
  \FunctionTok{cat}\NormalTok{(}\StringTok{"Número de titulaciones:"}\NormalTok{, }\FunctionTok{nrow}\NormalTok{(q1\_data), }\StringTok{"}\SpecialCharTok{\textbackslash{}n}\StringTok{"}\NormalTok{)}
  \FunctionTok{cat}\NormalTok{(}\StringTok{"Rendimiento promedio:"}\NormalTok{, scales}\SpecialCharTok{::}\FunctionTok{percent}\NormalTok{(}\FunctionTok{mean}\NormalTok{(q1\_data}\SpecialCharTok{$}\NormalTok{Tasa\_Rendimiento)), }\StringTok{"}\SpecialCharTok{\textbackslash{}n}\StringTok{"}\NormalTok{)}
  \FunctionTok{cat}\NormalTok{(}\StringTok{"Abandono promedio:"}\NormalTok{, scales}\SpecialCharTok{::}\FunctionTok{percent}\NormalTok{(}\FunctionTok{mean}\NormalTok{(q1\_data}\SpecialCharTok{$}\NormalTok{Tasa\_Abandono)), }\StringTok{"}\SpecialCharTok{\textbackslash{}n\textbackslash{}n}\StringTok{"}\NormalTok{)}
  \FunctionTok{cat}\NormalTok{(}\StringTok{"RECOMENDACIONES:}\SpecialCharTok{\textbackslash{}n}\StringTok{"}\NormalTok{)}
  \FunctionTok{cat}\NormalTok{(}\StringTok{"  {-} Mantener modelo actual de enseñanza}\SpecialCharTok{\textbackslash{}n}\StringTok{"}\NormalTok{)}
  \FunctionTok{cat}\NormalTok{(}\StringTok{"  {-} Documentar y replicar buenas prácticas en otras titulaciones}\SpecialCharTok{\textbackslash{}n}\StringTok{"}\NormalTok{)}
  \FunctionTok{cat}\NormalTok{(}\StringTok{"  {-} Usar como referencia para programas de mejora}\SpecialCharTok{\textbackslash{}n}\StringTok{"}\NormalTok{)}
  \FunctionTok{cat}\NormalTok{(}\StringTok{"  {-} Investigar factores de éxito (recursos, profesorado, estructura)}\SpecialCharTok{\textbackslash{}n}\StringTok{"}\NormalTok{)}
  \FunctionTok{cat}\NormalTok{(}\StringTok{"Titulaciones:}\SpecialCharTok{\textbackslash{}n}\StringTok{"}\NormalTok{)}
  \ControlFlowTok{for}\NormalTok{ (i }\ControlFlowTok{in} \DecValTok{1}\SpecialCharTok{:}\FunctionTok{nrow}\NormalTok{(q1\_data)) \{}
    \FunctionTok{cat}\NormalTok{(}\FunctionTok{sprintf}\NormalTok{(}\StringTok{"    • \%s (Rend: \%.1f\%\%, Aband: \%.1f\%\%)}\SpecialCharTok{\textbackslash{}n}\StringTok{"}\NormalTok{,}
\NormalTok{                q1\_data}\SpecialCharTok{$}\NormalTok{Titulacion[i],}
\NormalTok{                q1\_data}\SpecialCharTok{$}\NormalTok{Tasa\_Rendimiento[i] }\SpecialCharTok{*} \DecValTok{100}\NormalTok{,}
\NormalTok{                q1\_data}\SpecialCharTok{$}\NormalTok{Tasa\_Abandono[i] }\SpecialCharTok{*} \DecValTok{100}\NormalTok{))}
\NormalTok{  \}}
  \FunctionTok{cat}\NormalTok{(}\StringTok{"}\SpecialCharTok{\textbackslash{}n\textbackslash{}n}\StringTok{"}\NormalTok{)}
\NormalTok{\}}
\end{Highlighting}
\end{Shaded}

\begin{verbatim}
## Q1: TITULACIONES EXITOSAS | (Alto Rendimiento, Bajo Abandono)
## ---------------------------------------------------------------------------------------------------------------------------------- 
## Número de titulaciones: 76 
## Rendimiento promedio: 86% 
## Abandono promedio: 7% 
## 
## RECOMENDACIONES:
##   - Mantener modelo actual de enseñanza
##   - Documentar y replicar buenas prácticas en otras titulaciones
##   - Usar como referencia para programas de mejora
##   - Investigar factores de éxito (recursos, profesorado, estructura)
## Titulaciones:
##     • Grado en Bioquímica por la U.S. y UMA (Rend: 86.4%, Aband: 10.4%)
##     • Grado en Bioquímica por la Universidad de Sevilla y Universidad de Málaga (Rend: 83.9%, Aband: 11.3%)
##     • Grado en Ciencias de la Actividad Física y del Deporte (Rend: 88.1%, Aband: 10.5%)
##     • Grado en Comunicación Audiovisual (Rend: 88.8%, Aband: 11.6%)
##     • Grado en Criminología (Rend: 93.7%, Aband: 11.4%)
##     • Grado en Educación Infantil (Rend: 97.4%, Aband: 5.4%)
##     • Grado en Educación Primaria (Rend: 89.5%, Aband: 10.2%)
##     • Grado en Pedagogía (Rend: 91.5%, Aband: 12.6%)
##     • Grado en Publicidad y Relaciones Públicas (Rend: 89.7%, Aband: 12.0%)
##     • Grado en Óptica y Optometría (Rend: 88.7%, Aband: 12.5%)
##     • M. U. en Estudios Históricos Avanzados (Rend: 85.7%, Aband: 12.1%)
##     • M. U. en Gestión del Territorio, Instrumentos y Técnicas de Intervenc. (Rend: 82.4%, Aband: 5.6%)
##     • M.U. Profesorado de E.S.O. y Bachillerato, F. P. y E.de Idiomas (Rend: 98.6%, Aband: 2.9%)
##     • M.U. en Arte: Idea y Producción (Rend: 79.5%, Aband: 7.7%)
##     • M.U. en Artes del Espectáculo Vivo (R.D.1393/07) (Rend: 85.4%, Aband: 5.6%)
##     • M.U. en Asesoría Jurídico-Mercantil, Fiscal y Laboral (Rend: 80.7%, Aband: 5.5%)
##     • M.U. en Biología Avanzada: Investigación y Aplicación (Rend: 88.6%, Aband: 8.4%)
##     • M.U. en Ciencia y Tecnología de Nuevos Materiales (Rend: 88.2%, Aband: 12.7%)
##     • M.U. en Ciencia y Tecnología de Nuevos Materiales (R.D.1393/07) (Rend: 93.9%, Aband: 10.5%)
##     • M.U. en Ciudad y Arquitectura Sostenibles (Rend: 81.0%, Aband: 8.8%)
##     • M.U. en Comunicación Institucional y Política (Rend: 84.2%, Aband: 11.2%)
##     • M.U. en Comunicación Institucional y Política (R.D.1393/2007) (Rend: 78.8%, Aband: 6.0%)
##     • M.U. en Comunicación y Cultura (Rend: 78.9%, Aband: 6.0%)
##     • M.U. en Consultoría Laboral (Rend: 82.9%, Aband: 5.2%)
##     • M.U. en Derecho Constitucional (Rend: 78.3%, Aband: 10.6%)
##     • M.U. en Dirección y Planificación del Turismo (Rend: 80.0%, Aband: 11.0%)
##     • M.U. en Dirección y Planificación del Turismo (R.D.1393/07) (Rend: 79.2%, Aband: 5.4%)
##     • M.U. en Dirección, Evaluación y Calidad de las Instituc. de For (Rend: 88.8%, Aband: 3.7%)
##     • M.U. en Documentos y Libros. Archivos y Bibliotecas (Rend: 82.8%, Aband: 1.5%)
##     • M.U. en En.del Español como L. Extranjera y de Otras L. Modernas (Rend: 91.8%, Aband: 9.5%)
##     • M.U. en Especialización Profesional en Farmacia (Rend: 82.8%, Aband: 7.5%)
##     • M.U. en Estudios Americanos (Rend: 86.7%, Aband: 8.2%)
##     • M.U. en Estudios Americanos (R.D.1393/2007) (Rend: 88.4%, Aband: 9.8%)
##     • M.U. en Estudios Avanzados en Química (Rend: 88.9%, Aband: 3.3%)
##     • M.U. en Estudios Avanzados en Química (R.D.1393/07) (Rend: 94.0%, Aband: 0.0%)
##     • M.U. en Estudios Hispánicos Superiores (R.D.1393/07) (Rend: 113.3%, Aband: 7.2%)
##     • M.U. en Estudios Lingüísticos, Literarios y Culturales (Rend: 79.6%, Aband: 8.4%)
##     • M.U. en Estudios Lingüísticos, Literarios y Culturales por la U.S. (Rend: 85.8%, Aband: 7.0%)
##     • M.U. en Estudios de Género y Desarrollo Profesional (Rend: 85.7%, Aband: 11.8%)
##     • M.U. en Estudios de Género y Desarrollo Profesional (RD.07) (Rend: 80.8%, Aband: 9.2%)
##     • M.U. en Filosofía y Cultura Moderna (R.D.1393/2007) (Rend: 103.9%, Aband: 10.3%)
##     • M.U. en Formación y Orientación para el Trabajo (Rend: 88.6%, Aband: 1.4%)
##     • M.U. en Gestión Integral de la Edificación (Rend: 80.9%, Aband: 8.3%)
##     • M.U. en Gestión Integral de la Edificación (R.D.1393/07) (Rend: 81.4%, Aband: 7.1%)
##     • M.U. en Gestión del Territorio. Instrumentos y Técnicas de Intervenc (Rend: 78.9%, Aband: 10.6%)
##     • M.U. en Migraciones Internacionales, Salud y Bienestar: Modelos y Estrategias de Intervención (Rend: 83.7%, Aband: 11.0%)
##     • M.U. en Nuevas Tendencias Asistenciales en CC. de la Salud (Rend: 78.7%, Aband: 5.3%)
##     • M.U. en Nuevas Tendencias Asistenciales en Ciencias de la Salud (Rend: 77.6%, Aband: 8.4%)
##     • M.U. en Patrimonio Artístico Andaluz y su Proy. Iberoamericana (Rend: 83.2%, Aband: 6.7%)
##     • M.U. en Patrimonio Artístico Andaluz y su Proyección Iberoameric. (Rend: 80.0%, Aband: 10.8%)
##     • M.U. en Profesorado de E.S.O y Bachillerato, FP y E.Idiomas (Rend: 95.6%, Aband: 3.3%)
##     • M.U. en Psicología de la Intervención Social y Comunitaria (Rend: 86.1%, Aband: 3.4%)
##     • M.U. en Psicopedagogía (Rend: 92.8%, Aband: 2.7%)
##     • M.U. en Relaciones Jurídico-Privadas (R.D.1393/07) (Rend: 82.7%, Aband: 3.6%)
##     • M.U. en Seguridad Integral en Edificación (Rend: 88.1%, Aband: 12.7%)
##     • M.U. en Seguridad Integral en Edificación (R.D.1393/07) (Rend: 88.2%, Aband: 0.0%)
##     • M.U. en Seguridad Integral en la Industria y Prev. RR.LL. (Rend: 81.0%, Aband: 7.5%)
##     • M.U.en Enseñanza del Español como Lengua Extranjera y Otras Leng.Mod. (Rend: 86.1%, Aband: 5.6%)
##     • M.U.en Migraciones Internacionales, Salud y Bienestar: Mod. y Est. de Inter (Rend: 83.3%, Aband: 2.8%)
##     • M.U.en Seguridad Integral en la Industria y Prevención Riesgos Laborales (Rend: 84.5%, Aband: 9.2%)
##     • Máster Univ. en Psicología de la Intervención Social y Comunitaria (Rend: 89.9%, Aband: 4.5%)
##     • Máster Univers. en Profesorado de E.S.O y Bachillerato, FP y E.Idiomas (Rend: 94.0%, Aband: 3.0%)
##     • Máster Universitario en Biología Avanzada: Investigación y Aplicación (Rend: 87.9%, Aband: 8.1%)
##     • Máster Universitario en Ciencia y Tecnología de Nuevos Materiales (Rend: 85.8%, Aband: 10.4%)
##     • Máster Universitario en Documentos y Libros. Archivos y Bibliotecas (Rend: 79.2%, Aband: 4.1%)
##     • Máster Universitario en Estudios Americanos (Rend: 81.0%, Aband: 6.4%)
##     • Máster Universitario en Estudios Avanzados en Química (Rend: 84.3%, Aband: 5.6%)
##     • Máster Universitario en Estudios Lingüísticos, Literarios y Culturales (Rend: 77.4%, Aband: 8.3%)
##     • Máster Universitario en Estudios de Género y Desarrollo Profesional (Rend: 81.3%, Aband: 10.0%)
##     • Máster Universitario en Filosofía y Cultura Moderna (Rend: 84.0%, Aband: 9.8%)
##     • Máster Universitario en Formación y Orientación para el Trabajo (Rend: 87.3%, Aband: 3.6%)
##     • Máster Universitario en Psicopedagogía (Rend: 88.5%, Aband: 3.5%)
##     • Máster Unv. en Dirección, Evaluación y Calidad de las Instituc. de For (Rend: 84.5%, Aband: 6.9%)
##     • Máster Unv. en Investigación Médica: Clínica y Experimental (Rend: 81.6%, Aband: 8.8%)
##     • Máster Unv.en Patrimonio Artístico Andaluz y su Proyección Iberoameric. (Rend: 84.4%, Aband: 9.7%)
##     • Máster Unv.en Patrimonio Artístico Andaluz y su Proyección Iberoamericana (Rend: 77.5%, Aband: 5.1%)
\end{verbatim}

\begin{Shaded}
\begin{Highlighting}[]
\CommentTok{\# Q2: ABANDONO SELECTIVO}
\NormalTok{q2\_data }\OtherTok{\textless{}{-}}\NormalTok{ titulaciones\_quadrants }\SpecialCharTok{\%\textgreater{}\%} \FunctionTok{filter}\NormalTok{(}\FunctionTok{str\_detect}\NormalTok{(Cuadrante, }\StringTok{"Q2"}\NormalTok{))}
\ControlFlowTok{if}\NormalTok{ (}\FunctionTok{nrow}\NormalTok{(q2\_data) }\SpecialCharTok{\textgreater{}} \DecValTok{0}\NormalTok{) \{}
  \FunctionTok{cat}\NormalTok{(}\StringTok{"Q2: ABANDONO SELECTIVO | (Alto Rendimiento, Alto Abandono)}\SpecialCharTok{\textbackslash{}n}\StringTok{"}\NormalTok{)}
  \FunctionTok{cat}\NormalTok{(}\FunctionTok{strrep}\NormalTok{(}\StringTok{"{-}"}\NormalTok{, }\DecValTok{130}\NormalTok{), }\StringTok{"}\SpecialCharTok{\textbackslash{}n}\StringTok{"}\NormalTok{)}
  \FunctionTok{cat}\NormalTok{(}\StringTok{"Número de titulaciones:"}\NormalTok{, }\FunctionTok{nrow}\NormalTok{(q2\_data), }\StringTok{"}\SpecialCharTok{\textbackslash{}n}\StringTok{"}\NormalTok{)}
  \FunctionTok{cat}\NormalTok{(}\StringTok{"Rendimiento promedio:"}\NormalTok{, scales}\SpecialCharTok{::}\FunctionTok{percent}\NormalTok{(}\FunctionTok{mean}\NormalTok{(q2\_data}\SpecialCharTok{$}\NormalTok{Tasa\_Rendimiento)), }\StringTok{"}\SpecialCharTok{\textbackslash{}n}\StringTok{"}\NormalTok{)}
  \FunctionTok{cat}\NormalTok{(}\StringTok{"Abandono promedio:"}\NormalTok{, scales}\SpecialCharTok{::}\FunctionTok{percent}\NormalTok{(}\FunctionTok{mean}\NormalTok{(q2\_data}\SpecialCharTok{$}\NormalTok{Tasa\_Abandono)), }\StringTok{"}\SpecialCharTok{\textbackslash{}n\textbackslash{}n}\StringTok{"}\NormalTok{)}
  \FunctionTok{cat}\NormalTok{(}\StringTok{"RECOMENDACIONES:}\SpecialCharTok{\textbackslash{}n}\StringTok{"}\NormalTok{)}
  \FunctionTok{cat}\NormalTok{(}\StringTok{"  {-} Investigar causas del abandono (académicas vs no académicas)}\SpecialCharTok{\textbackslash{}n}\StringTok{"}\NormalTok{)}
  \FunctionTok{cat}\NormalTok{(}\StringTok{"  {-} Análisis detallado: ¿abandonan estudiantes con bajo rendimiento o con bueno?}\SpecialCharTok{\textbackslash{}n}\StringTok{"}\NormalTok{)}
  \FunctionTok{cat}\NormalTok{(}\StringTok{"  {-} Posibles causas no académicas: horarios incompatibles, coste, compatibilidad laboral}\SpecialCharTok{\textbackslash{}n}\StringTok{"}\NormalTok{)}
  \FunctionTok{cat}\NormalTok{(}\StringTok{"  {-} Considerar modalidades alternativas (semipresencial, flexible)}\SpecialCharTok{\textbackslash{}n}\StringTok{"}\NormalTok{)}
  \FunctionTok{cat}\NormalTok{(}\StringTok{"Titulaciones:}\SpecialCharTok{\textbackslash{}n}\StringTok{"}\NormalTok{)}
  \ControlFlowTok{for}\NormalTok{ (i }\ControlFlowTok{in} \DecValTok{1}\SpecialCharTok{:}\FunctionTok{nrow}\NormalTok{(q2\_data)) \{}
    \FunctionTok{cat}\NormalTok{(}\FunctionTok{sprintf}\NormalTok{(}\StringTok{"    • \%s (Rend: \%.1f\%\%, Aband: \%.1f\%\%)}\SpecialCharTok{\textbackslash{}n}\StringTok{"}\NormalTok{,}
\NormalTok{                q2\_data}\SpecialCharTok{$}\NormalTok{Titulacion[i],}
\NormalTok{                q2\_data}\SpecialCharTok{$}\NormalTok{Tasa\_Rendimiento[i] }\SpecialCharTok{*} \DecValTok{100}\NormalTok{,}
\NormalTok{                q2\_data}\SpecialCharTok{$}\NormalTok{Tasa\_Abandono[i] }\SpecialCharTok{*} \DecValTok{100}\NormalTok{))}
\NormalTok{  \}}
  \FunctionTok{cat}\NormalTok{(}\StringTok{"}\SpecialCharTok{\textbackslash{}n\textbackslash{}n}\StringTok{"}\NormalTok{)}
\NormalTok{\}}
\end{Highlighting}
\end{Shaded}

\begin{verbatim}
## Q2: ABANDONO SELECTIVO | (Alto Rendimiento, Alto Abandono)
## ---------------------------------------------------------------------------------------------------------------------------------- 
## Número de titulaciones: 23 
## Rendimiento promedio: 83% 
## Abandono promedio: 18% 
## 
## RECOMENDACIONES:
##   - Investigar causas del abandono (académicas vs no académicas)
##   - Análisis detallado: ¿abandonan estudiantes con bajo rendimiento o con bueno?
##   - Posibles causas no académicas: horarios incompatibles, coste, compatibilidad laboral
##   - Considerar modalidades alternativas (semipresencial, flexible)
## Titulaciones:
##     • Grado en Bellas Artes (Rend: 86.9%, Aband: 18.2%)
##     • Grado en Biomedicina Básica y Experimental (Rend: 91.8%, Aband: 17.1%)
##     • Grado en Bioquímica por la U. de Sevilla y U. de Málaga (Rend: 86.9%, Aband: 13.6%)
##     • Grado en Conservación y Restauración de Bienes Culturales (Rend: 84.3%, Aband: 25.9%)
##     • Grado en Enfermería (Rend: 86.4%, Aband: 21.5%)
##     • Grado en Estudios de Asia Oriental por la U.S. y la U.de Málaga (Rend: 84.0%, Aband: 21.2%)
##     • Grado en Estudios de Asia Oriental por la US y UMA (Rend: 86.1%, Aband: 17.3%)
##     • Grado en Estudios de Asia Oriental por la Unv.de Sevilla y la Unv.de Málaga (Rend: 82.9%, Aband: 17.8%)
##     • Grado en Fisioterapia (Rend: 86.7%, Aband: 26.1%)
##     • Grado en Periodismo (Rend: 79.3%, Aband: 14.8%)
##     • Grado en Podología (Rend: 79.3%, Aband: 28.2%)
##     • Grado en Psicología (Rend: 86.9%, Aband: 12.9%)
##     • M.U. en Comunicación y Cultura (R.D.1393/07) (Rend: 85.0%, Aband: 13.8%)
##     • M.U. en Derecho Constitucional (R.D.1393/07) (Rend: 82.1%, Aband: 19.2%)
##     • M.U. en Dirección, Evaluación y Calidad de las Instituciones de For. (Rend: 86.8%, Aband: 15.2%)
##     • M.U. en Escritura Creativa (Rend: 78.6%, Aband: 12.9%)
##     • M.U. en Escritura Creativa (R.D.1393/07) (Rend: 76.6%, Aband: 14.1%)
##     • M.U. en Estudios Históricos Avanzados (Rend: 84.2%, Aband: 15.5%)
##     • M.U. en Investigación Médica: Clínica y Experimental (Rend: 81.0%, Aband: 14.1%)
##     • Máster Un. en Asesoría Jurídico-Mercantil, Fiscal y Laboral (Rend: 79.3%, Aband: 30.8%)
##     • Máster Universitario en Consultoría Laboral (Rend: 80.5%, Aband: 14.4%)
##     • Máster Universitario en Estudios Históricos Avanzados (Rend: 82.8%, Aband: 14.0%)
##     • Máster Universitario en Relaciones Jurídico-Privadas (Rend: 77.9%, Aband: 23.7%)
\end{verbatim}

\begin{Shaded}
\begin{Highlighting}[]
\CommentTok{\# Q3: PROBLEMÁTICO}
\NormalTok{q3\_data }\OtherTok{\textless{}{-}}\NormalTok{ titulaciones\_quadrants }\SpecialCharTok{\%\textgreater{}\%} \FunctionTok{filter}\NormalTok{(}\FunctionTok{str\_detect}\NormalTok{(Cuadrante, }\StringTok{"Q3"}\NormalTok{))}
\ControlFlowTok{if}\NormalTok{ (}\FunctionTok{nrow}\NormalTok{(q3\_data) }\SpecialCharTok{\textgreater{}} \DecValTok{0}\NormalTok{) \{}
  \FunctionTok{cat}\NormalTok{(}\StringTok{"Q3: PROBLEMÁTICO | (Bajo Rendimiento, Bajo Abandono)}\SpecialCharTok{\textbackslash{}n}\StringTok{"}\NormalTok{)}
  \FunctionTok{cat}\NormalTok{(}\FunctionTok{strrep}\NormalTok{(}\StringTok{"{-}"}\NormalTok{, }\DecValTok{130}\NormalTok{), }\StringTok{"}\SpecialCharTok{\textbackslash{}n}\StringTok{"}\NormalTok{)}
  \FunctionTok{cat}\NormalTok{(}\StringTok{"Número de titulaciones:"}\NormalTok{, }\FunctionTok{nrow}\NormalTok{(q3\_data), }\StringTok{"}\SpecialCharTok{\textbackslash{}n}\StringTok{"}\NormalTok{)}
  \FunctionTok{cat}\NormalTok{(}\StringTok{"Rendimiento promedio:"}\NormalTok{, scales}\SpecialCharTok{::}\FunctionTok{percent}\NormalTok{(}\FunctionTok{mean}\NormalTok{(q3\_data}\SpecialCharTok{$}\NormalTok{Tasa\_Rendimiento)), }\StringTok{"}\SpecialCharTok{\textbackslash{}n}\StringTok{"}\NormalTok{)}
  \FunctionTok{cat}\NormalTok{(}\StringTok{"Abandono promedio:"}\NormalTok{, scales}\SpecialCharTok{::}\FunctionTok{percent}\NormalTok{(}\FunctionTok{mean}\NormalTok{(q3\_data}\SpecialCharTok{$}\NormalTok{Tasa\_Abandono)), }\StringTok{"}\SpecialCharTok{\textbackslash{}n\textbackslash{}n}\StringTok{"}\NormalTok{)}
  \FunctionTok{cat}\NormalTok{(}\StringTok{"RECOMENDACIONES:}\SpecialCharTok{\textbackslash{}n}\StringTok{"}\NormalTok{)}
  \FunctionTok{cat}\NormalTok{(}\StringTok{"  {-} Intervención moderada prioritaria}\SpecialCharTok{\textbackslash{}n}\StringTok{"}\NormalTok{)}
  \FunctionTok{cat}\NormalTok{(}\StringTok{"  {-} Estudiantes se quedan pero arrastrando asignaturas (acumulan suspensos)}\SpecialCharTok{\textbackslash{}n}\StringTok{"}\NormalTok{)}
  \FunctionTok{cat}\NormalTok{(}\StringTok{"  {-} Revisar plan de estudios y secuenciación de asignaturas}\SpecialCharTok{\textbackslash{}n}\StringTok{"}\NormalTok{)}
  \FunctionTok{cat}\NormalTok{(}\StringTok{"  {-} Fortalecer apoyo académico (tutorías, refuerzo en asignaturas básicas)}\SpecialCharTok{\textbackslash{}n}\StringTok{"}\NormalTok{)}
  \FunctionTok{cat}\NormalTok{(}\StringTok{"Titulaciones:}\SpecialCharTok{\textbackslash{}n}\StringTok{"}\NormalTok{)}
  \ControlFlowTok{for}\NormalTok{ (i }\ControlFlowTok{in} \DecValTok{1}\SpecialCharTok{:}\FunctionTok{nrow}\NormalTok{(q3\_data)) \{}
    \FunctionTok{cat}\NormalTok{(}\FunctionTok{sprintf}\NormalTok{(}\StringTok{"    • \%s (Rend: \%.1f\%\%, Aband: \%.1f\%\%)}\SpecialCharTok{\textbackslash{}n}\StringTok{"}\NormalTok{,}
\NormalTok{                q3\_data}\SpecialCharTok{$}\NormalTok{Titulacion[i],}
\NormalTok{                q3\_data}\SpecialCharTok{$}\NormalTok{Tasa\_Rendimiento[i] }\SpecialCharTok{*} \DecValTok{100}\NormalTok{,}
\NormalTok{                q3\_data}\SpecialCharTok{$}\NormalTok{Tasa\_Abandono[i] }\SpecialCharTok{*} \DecValTok{100}\NormalTok{))}
\NormalTok{  \}}
  \FunctionTok{cat}\NormalTok{(}\StringTok{"}\SpecialCharTok{\textbackslash{}n\textbackslash{}n}\StringTok{"}\NormalTok{)}
\NormalTok{\}}
\end{Highlighting}
\end{Shaded}

\begin{verbatim}
## Q3: PROBLEMÁTICO | (Bajo Rendimiento, Bajo Abandono)
## ---------------------------------------------------------------------------------------------------------------------------------- 
## Número de titulaciones: 23 
## Rendimiento promedio: 69% 
## Abandono promedio: 9% 
## 
## RECOMENDACIONES:
##   - Intervención moderada prioritaria
##   - Estudiantes se quedan pero arrastrando asignaturas (acumulan suspensos)
##   - Revisar plan de estudios y secuenciación de asignaturas
##   - Fortalecer apoyo académico (tutorías, refuerzo en asignaturas básicas)
## Titulaciones:
##     • M.U. en Arqueología por la US y UGR (Rend: 73.7%, Aband: 10.5%)
##     • M.U. en Arquitectura (Rend: 72.4%, Aband: 4.2%)
##     • M.U. en Arte: Idea y Producción (R.D.1393/07) (Rend: 76.3%, Aband: 7.7%)
##     • M.U. en Filosofía y Cultura Moderna (Rend: 70.9%, Aband: 11.1%)
##     • M.U. en Física Nuclear (Rend: 71.1%, Aband: 8.9%)
##     • M.U. en Física Nuclear (Máster Interuniversitario) (R.D.1393/07) (Rend: 68.9%, Aband: 6.2%)
##     • M.U. en Gestión Estratégica y Negocios Internacionales (Rend: 76.3%, Aband: 10.9%)
##     • M.U. en Matemáticas (Rend: 72.3%, Aband: 11.9%)
##     • M.U. en Organización Industrial y Gestión de Empresas (Rend: 64.7%, Aband: 12.9%)
##     • M.U. en Peritación y Reparación de Edificios (Rend: 68.6%, Aband: 4.3%)
##     • M.U. en Sistemas de Energía Térmica (Rend: 65.6%, Aband: 10.6%)
##     • Máster Univ.en Nuevas Tendencias Asistenciales en CC. de la Salud (Rend: 59.8%, Aband: 7.6%)
##     • Máster Universitario en Arqueología por la U.de Sevilla y U.de Granada (Rend: 63.2%, Aband: 10.6%)
##     • Máster Universitario en Arquitectura (Rend: 51.3%, Aband: 6.4%)
##     • Máster Universitario en Arte: Idea y Producción (Rend: 73.2%, Aband: 4.7%)
##     • Máster Universitario en Artes del Espectáculo Vivo (Rend: 71.6%, Aband: 12.2%)
##     • Máster Universitario en Asesoría Jurídico-Mercantil, Fiscal y Laboral (Rend: 66.2%, Aband: 6.2%)
##     • Máster Universitario en Ciudad y Arquitectura Sostenibles (Rend: 76.3%, Aband: 6.8%)
##     • Máster Universitario en Escritura Creativa (Rend: 73.9%, Aband: 5.6%)
##     • Máster Universitario en Física Nuclear por la UAM,UCM,UB,UGR,USAL y la US (Rend: 71.3%, Aband: 9.9%)
##     • Máster Universitario en Lógica, Computación e Inteligencia Artificial (Rend: 73.5%, Aband: 12.5%)
##     • Máster Universitario en Matemáticas (Rend: 70.9%, Aband: 10.6%)
##     • Máster Universitario en Organización Industrial y Gestión de Empresas (Rend: 56.4%, Aband: 10.4%)
\end{verbatim}

\begin{Shaded}
\begin{Highlighting}[]
\CommentTok{\# Q4: EN RIESGO}
\NormalTok{q4\_data }\OtherTok{\textless{}{-}}\NormalTok{ titulaciones\_quadrants }\SpecialCharTok{\%\textgreater{}\%} \FunctionTok{filter}\NormalTok{(}\FunctionTok{str\_detect}\NormalTok{(Cuadrante, }\StringTok{"Q4"}\NormalTok{))}
\ControlFlowTok{if}\NormalTok{ (}\FunctionTok{nrow}\NormalTok{(q4\_data) }\SpecialCharTok{\textgreater{}} \DecValTok{0}\NormalTok{) \{}
  \FunctionTok{cat}\NormalTok{(}\StringTok{"Q4: EN RIESGO | (Bajo Rendimiento, Alto Abandono) PRIORIDAD MÁXIMA}\SpecialCharTok{\textbackslash{}n}\StringTok{"}\NormalTok{)}
  \FunctionTok{cat}\NormalTok{(}\FunctionTok{strrep}\NormalTok{(}\StringTok{"{-}"}\NormalTok{, }\DecValTok{130}\NormalTok{), }\StringTok{"}\SpecialCharTok{\textbackslash{}n}\StringTok{"}\NormalTok{)}
  \FunctionTok{cat}\NormalTok{(}\StringTok{"Número de titulaciones:"}\NormalTok{, }\FunctionTok{nrow}\NormalTok{(q4\_data), }\StringTok{"}\SpecialCharTok{\textbackslash{}n}\StringTok{"}\NormalTok{)}
  \FunctionTok{cat}\NormalTok{(}\StringTok{"Rendimiento promedio:"}\NormalTok{, scales}\SpecialCharTok{::}\FunctionTok{percent}\NormalTok{(}\FunctionTok{mean}\NormalTok{(q4\_data}\SpecialCharTok{$}\NormalTok{Tasa\_Rendimiento)), }\StringTok{"}\SpecialCharTok{\textbackslash{}n}\StringTok{"}\NormalTok{)}
  \FunctionTok{cat}\NormalTok{(}\StringTok{"Abandono promedio:"}\NormalTok{, scales}\SpecialCharTok{::}\FunctionTok{percent}\NormalTok{(}\FunctionTok{mean}\NormalTok{(q4\_data}\SpecialCharTok{$}\NormalTok{Tasa\_Abandono)), }\StringTok{"}\SpecialCharTok{\textbackslash{}n\textbackslash{}n}\StringTok{"}\NormalTok{)}
  \FunctionTok{cat}\NormalTok{(}\StringTok{"RECOMENDACIONES: (INTERVENCIÓN URGENTE E INMEDIATA)}\SpecialCharTok{\textbackslash{}n}\StringTok{"}\NormalTok{)}
  \FunctionTok{cat}\NormalTok{(}\StringTok{"  {-} Posibles problemas estructurales o expectativas no cumplidas}\SpecialCharTok{\textbackslash{}n}\StringTok{"}\NormalTok{)}
  \FunctionTok{cat}\NormalTok{(}\StringTok{"  {-} Revisión profunda del plan de estudios (objetivos, contenidos, metodología)}\SpecialCharTok{\textbackslash{}n}\StringTok{"}\NormalTok{)}
  \FunctionTok{cat}\NormalTok{(}\StringTok{"  {-} Evaluación completa del cuerpo docente y recursos disponibles}\SpecialCharTok{\textbackslash{}n}\StringTok{"}\NormalTok{)}
  \FunctionTok{cat}\NormalTok{(}\StringTok{"  {-} Mejorar información pre{-}matriculación (expectativas realistas)}\SpecialCharTok{\textbackslash{}n}\StringTok{"}\NormalTok{)}
  \FunctionTok{cat}\NormalTok{(}\StringTok{"Titulaciones:}\SpecialCharTok{\textbackslash{}n}\StringTok{"}\NormalTok{)}
  \ControlFlowTok{for}\NormalTok{ (i }\ControlFlowTok{in} \DecValTok{1}\SpecialCharTok{:}\FunctionTok{nrow}\NormalTok{(q4\_data)) \{}
    \FunctionTok{cat}\NormalTok{(}\FunctionTok{sprintf}\NormalTok{(}\StringTok{"    • \%s (Rend: \%.1f\%\%, Aband: \%.1f\%\%)}\SpecialCharTok{\textbackslash{}n}\StringTok{"}\NormalTok{,}
\NormalTok{                q4\_data}\SpecialCharTok{$}\NormalTok{Titulacion[i],}
\NormalTok{                q4\_data}\SpecialCharTok{$}\NormalTok{Tasa\_Rendimiento[i] }\SpecialCharTok{*} \DecValTok{100}\NormalTok{,}
\NormalTok{                q4\_data}\SpecialCharTok{$}\NormalTok{Tasa\_Abandono[i] }\SpecialCharTok{*} \DecValTok{100}\NormalTok{))}
\NormalTok{  \}}
  \FunctionTok{cat}\NormalTok{(}\StringTok{"}\SpecialCharTok{\textbackslash{}n\textbackslash{}n}\StringTok{"}\NormalTok{)}
\NormalTok{\}}
\end{Highlighting}
\end{Shaded}

\begin{verbatim}
## Q4: EN RIESGO | (Bajo Rendimiento, Alto Abandono) PRIORIDAD MÁXIMA
## ---------------------------------------------------------------------------------------------------------------------------------- 
## Número de titulaciones: 76 
## Rendimiento promedio: 62% 
## Abandono promedio: 28% 
## 
## RECOMENDACIONES: (INTERVENCIÓN URGENTE E INMEDIATA)
##   - Posibles problemas estructurales o expectativas no cumplidas
##   - Revisión profunda del plan de estudios (objetivos, contenidos, metodología)
##   - Evaluación completa del cuerpo docente y recursos disponibles
##   - Mejorar información pre-matriculación (expectativas realistas)
## Titulaciones:
##     • Grado en Administración y Dirección de Empresas (Rend: 57.9%, Aband: 37.1%)
##     • Grado en Antropología Social y Cultural (Rend: 62.5%, Aband: 45.2%)
##     • Grado en Biología (Rend: 76.1%, Aband: 15.7%)
##     • Grado en Derecho (Rend: 72.1%, Aband: 14.6%)
##     • Grado en Economía (Rend: 60.5%, Aband: 32.0%)
##     • Grado en Estudios Franceses (Rend: 62.0%, Aband: 39.5%)
##     • Grado en Estudios Ingleses (Rend: 70.2%, Aband: 25.3%)
##     • Grado en Estudios Árabes e Islámicos (Rend: 65.4%, Aband: 43.0%)
##     • Grado en Filología Clásica (Rend: 67.2%, Aband: 55.9%)
##     • Grado en Filología Hispánica (Rend: 62.6%, Aband: 32.3%)
##     • Grado en Filosofía (Rend: 67.3%, Aband: 44.3%)
##     • Grado en Finanzas y Contabilidad (Rend: 48.9%, Aband: 49.3%)
##     • Grado en Física (Rend: 64.6%, Aband: 26.8%)
##     • Grado en Geografia y Gestión del Territorio (Rend: 66.4%, Aband: 33.2%)
##     • Grado en Geografía y Gestión del Territorio (Rend: 66.0%, Aband: 42.7%)
##     • Grado en Historia (Rend: 66.2%, Aband: 26.5%)
##     • Grado en Historia del Arte (Rend: 62.1%, Aband: 39.2%)
##     • Grado en Ingeniería Aeroespacial (Rend: 66.3%, Aband: 16.0%)
##     • Grado en Ingeniería Civil (Rend: 51.2%, Aband: 45.0%)
##     • Grado en Ingeniería Electrónica Industrial (Rend: 43.5%, Aband: 46.1%)
##     • Grado en Ingeniería Electrónica, Robótica y Mecatrónica (UMA-US) (Rend: 67.8%, Aband: 17.6%)
##     • Grado en Ingeniería Eléctrica (Rend: 46.3%, Aband: 44.0%)
##     • Grado en Ingeniería Informática-Ingeniería de Computadores (Rend: 49.7%, Aband: 59.2%)
##     • Grado en Ingeniería Informática-Ingeniería del Software (Rend: 70.7%, Aband: 26.3%)
##     • Grado en Ingeniería Informática-Tecnologías Informáticas (Rend: 56.4%, Aband: 45.3%)
##     • Grado en Ingeniería Mecánica (Rend: 51.7%, Aband: 30.9%)
##     • Grado en Ingeniería Química (Rend: 52.4%, Aband: 36.3%)
##     • Grado en Ingeniería Química Industrial (Rend: 45.6%, Aband: 44.5%)
##     • Grado en Ingeniería de Materiales (Rend: 64.5%, Aband: 23.5%)
##     • Grado en Ingeniería de Organización Industrial por la U. M. y U.S. (Rend: 65.9%, Aband: 24.1%)
##     • Grado en Ingeniería de Organización Industrial por la U. Málaga y U.Sevilla (Rend: 62.0%, Aband: 27.0%)
##     • Grado en Ingeniería de Organización Industrial por la UMA y US (Rend: 65.2%, Aband: 14.5%)
##     • Grado en Ingeniería de Tecnologías Industriales (Rend: 62.4%, Aband: 16.8%)
##     • Grado en Ingeniería de la Energía por la U.S y la U. Málaga (Rend: 56.0%, Aband: 23.6%)
##     • Grado en Ingeniería de la Energía por la US y UMA (Rend: 58.3%, Aband: 37.3%)
##     • Grado en Ingeniería de la Energía por la Un. de Sevilla y la Un. de Málaga (Rend: 55.1%, Aband: 28.1%)
##     • Grado en Ingeniería de la Salud por la U. Málaga y la U.Sevilla (Rend: 68.6%, Aband: 24.4%)
##     • Grado en Ingeniería de la Salud por la UMA y US (Rend: 68.9%, Aband: 17.4%)
##     • Grado en Ingeniería de la Salud por la Univ. de Málaga y la Univ.de Sevilla (Rend: 73.2%, Aband: 14.8%)
##     • Grado en Ingeniería de las Tecnologías de Telecomunicación (Rend: 53.1%, Aband: 32.7%)
##     • Grado en Ingeniería en Diseño Industrial y Desarrollo del Producto (Rend: 61.1%, Aband: 22.6%)
##     • Grado en Lengua y Literatura Alemanas (Rend: 66.6%, Aband: 37.0%)
##     • Grado en Marketing e Investigación de Mercados (Rend: 75.4%, Aband: 26.5%)
##     • Grado en Matemáticas (Rend: 66.7%, Aband: 23.2%)
##     • Grado en Química (Rend: 62.0%, Aband: 27.0%)
##     • Grado en Relaciones Laborales y Recursos Humanos (Rend: 65.1%, Aband: 22.3%)
##     • Grado en Turismo (Rend: 64.5%, Aband: 29.9%)
##     • M.U. en Antropología: Gestión de la Diver. Cul., el Patr. y el Des. (Rend: 73.4%, Aband: 17.6%)
##     • M.U. en Antropología: Gestión de la Diversidad Cultural, el Pat.y el Des. (Rend: 66.0%, Aband: 16.6%)
##     • M.U. en Arqueología (R.D.1393/07) (Rend: 71.4%, Aband: 16.3%)
##     • M.U. en Artes del Espectáculo Vivo (Rend: 75.7%, Aband: 17.1%)
##     • M.U. en Diseño Avanzado en Ingeniería Mecánica (Rend: 56.4%, Aband: 32.2%)
##     • M.U. en Economía y Desarrollo (Rend: 69.0%, Aband: 19.6%)
##     • M.U. en Economía y Desarrollo (R.D.1393/07) (Rend: 69.3%, Aband: 17.3%)
##     • M.U. en Estudios Hispánicos Superiores (Rend: 75.5%, Aband: 13.5%)
##     • M.U. en Lógica, Computación e Inteligencia Artificial (Rend: 62.5%, Aband: 18.0%)
##     • M.U. en Microelectrónica: Diseño y Aplicaciones de Sistemas (Rend: 52.5%, Aband: 31.9%)
##     • M.U. en Microelectrónica: Dño. y Apps.de Sistemas M./Nano. (Rend: 39.9%, Aband: 37.9%)
##     • M.U. en Relaciones Jurídicas-Privadas (Rend: 69.9%, Aband: 18.1%)
##     • M.U. en Sistemas de Energía Eléctrica (Rend: 54.0%, Aband: 21.7%)
##     • Máster Universitario en Comunicación Institucional y Política (Rend: 45.1%, Aband: 14.6%)
##     • Máster Universitario en Comunicación y Cultura (Rend: 51.1%, Aband: 16.3%)
##     • Máster Universitario en Derecho Constitucional (Rend: 73.0%, Aband: 27.9%)
##     • Máster Universitario en Dirección y Planificación del Turismo (Rend: 72.0%, Aband: 20.5%)
##     • Máster Universitario en Diseño Avanzado en Ingeniería Mecánica (Rend: 44.8%, Aband: 28.2%)
##     • Máster Universitario en Economía y Desarrollo (Rend: 56.5%, Aband: 25.9%)
##     • Máster Universitario en Especialización Profesional en Farmacia (Rend: 45.0%, Aband: 17.0%)
##     • Máster Universitario en Estudios Hispánicos Superiores (Rend: 70.8%, Aband: 28.9%)
##     • Máster Universitario en Gestión Estratégica y Negocios Internacionales (Rend: 75.2%, Aband: 15.1%)
##     • Máster Universitario en Gestión Integral de la Edificación (Rend: 73.4%, Aband: 16.0%)
##     • Máster Universitario en Peritación y Reparación de Edificios (Rend: 66.4%, Aband: 26.8%)
##     • Máster Universitario en Seguridad Integral en Edificación (Rend: 75.6%, Aband: 25.3%)
##     • Máster Universitario en Sistemas de Energía Eléctrica (Rend: 52.8%, Aband: 14.8%)
##     • Máster Universitario en Sistemas de Energía Térmica (Rend: 56.7%, Aband: 31.7%)
##     • Máster Unv. en Gestión del Territorio. Instrumentos y Técnicas de Intervenc (Rend: 69.1%, Aband: 18.3%)
##     • Máster Unv. en Microelectrónica: Diseño y Aplicaciones de Sistemas (Rend: 57.7%, Aband: 35.8%)
\end{verbatim}

\begin{center}\rule{0.5\linewidth}{0.5pt}\end{center}

\section{10. Conclusiones}\label{conclusiones}

\begin{Shaded}
\begin{Highlighting}[]
\FunctionTok{cat}\NormalTok{(}\StringTok{"CONCLUSIONES DEL ANÁLISIS}\SpecialCharTok{\textbackslash{}n}\StringTok{"}\NormalTok{)}
\end{Highlighting}
\end{Shaded}

\begin{verbatim}
## CONCLUSIONES DEL ANÁLISIS
\end{verbatim}

\begin{Shaded}
\begin{Highlighting}[]
\FunctionTok{cat}\NormalTok{(}\FunctionTok{strrep}\NormalTok{(}\StringTok{"="}\NormalTok{, }\DecValTok{100}\NormalTok{), }\StringTok{"}\SpecialCharTok{\textbackslash{}n\textbackslash{}n}\StringTok{"}\NormalTok{)}
\end{Highlighting}
\end{Shaded}

\begin{verbatim}
## ====================================================================================================
\end{verbatim}

\begin{Shaded}
\begin{Highlighting}[]
\FunctionTok{cat}\NormalTok{(}\StringTok{"1. VALIDACIÓN DE HIPÓTESIS:}\SpecialCharTok{\textbackslash{}n}\StringTok{"}\NormalTok{)}
\end{Highlighting}
\end{Shaded}

\begin{verbatim}
## 1. VALIDACIÓN DE HIPÓTESIS:
\end{verbatim}

\begin{Shaded}
\begin{Highlighting}[]
\FunctionTok{cat}\NormalTok{(}\StringTok{"   H₁ ACEPTADA: Las titulaciones pueden agruparse en clústeres definidos}\SpecialCharTok{\textbackslash{}n}\StringTok{"}\NormalTok{)}
\end{Highlighting}
\end{Shaded}

\begin{verbatim}
##    H₁ ACEPTADA: Las titulaciones pueden agruparse en clústeres definidos
\end{verbatim}

\begin{Shaded}
\begin{Highlighting}[]
\FunctionTok{cat}\NormalTok{(}\StringTok{"   que facilitan la intervención institucional.}\SpecialCharTok{\textbackslash{}n\textbackslash{}n}\StringTok{"}\NormalTok{)}
\end{Highlighting}
\end{Shaded}

\begin{verbatim}
##    que facilitan la intervención institucional.
\end{verbatim}

\begin{Shaded}
\begin{Highlighting}[]
\FunctionTok{cat}\NormalTok{(}\StringTok{"2. ESTRUCTURA DE CLÚSTERES IDENTIFICADA:}\SpecialCharTok{\textbackslash{}n}\StringTok{"}\NormalTok{)}
\end{Highlighting}
\end{Shaded}

\begin{verbatim}
## 2. ESTRUCTURA DE CLÚSTERES IDENTIFICADA:
\end{verbatim}

\begin{Shaded}
\begin{Highlighting}[]
\NormalTok{n\_clusters }\OtherTok{\textless{}{-}} \FunctionTok{nrow}\NormalTok{(cluster\_summary)}
\ControlFlowTok{for}\NormalTok{ (i }\ControlFlowTok{in} \DecValTok{1}\SpecialCharTok{:}\NormalTok{n\_clusters) \{}
\NormalTok{  row }\OtherTok{\textless{}{-}}\NormalTok{ cluster\_summary[i, ]}
  \FunctionTok{cat}\NormalTok{(}\FunctionTok{sprintf}\NormalTok{(}\StringTok{"   • Clúster \%d: \%d titulaciones}\SpecialCharTok{\textbackslash{}n}\StringTok{"}\NormalTok{, i, row}\SpecialCharTok{$}\NormalTok{n\_titulaciones))}
  \FunctionTok{cat}\NormalTok{(}\FunctionTok{sprintf}\NormalTok{(}\StringTok{"     {-} Rendimiento: \%.2f\%\% ± \%.2f\%\%}\SpecialCharTok{\textbackslash{}n}\StringTok{"}\NormalTok{, }
\NormalTok{              row}\SpecialCharTok{$}\NormalTok{Rendimiento\_Media }\SpecialCharTok{*} \DecValTok{100}\NormalTok{, row}\SpecialCharTok{$}\NormalTok{Rendimiento\_SD }\SpecialCharTok{*} \DecValTok{100}\NormalTok{))}
  \FunctionTok{cat}\NormalTok{(}\FunctionTok{sprintf}\NormalTok{(}\StringTok{"     {-} Abandono: \%.2f\%\% ± \%.2f\%\%}\SpecialCharTok{\textbackslash{}n\textbackslash{}n}\StringTok{"}\NormalTok{, }
\NormalTok{              row}\SpecialCharTok{$}\NormalTok{Abandono\_Media }\SpecialCharTok{*} \DecValTok{100}\NormalTok{, row}\SpecialCharTok{$}\NormalTok{Abandono\_SD }\SpecialCharTok{*} \DecValTok{100}\NormalTok{))}
\NormalTok{\}}
\end{Highlighting}
\end{Shaded}

\begin{verbatim}
##    • Clúster 1: 59 titulaciones
##      - Rendimiento: 59.27% ± 8.66%
##      - Abandono: 31.06% ± 10.84%
## 
##    • Clúster 2: 139 titulaciones
##      - Rendimiento: 81.08% ± 8.54%
##      - Abandono: 10.71% ± 5.62%
\end{verbatim}

\begin{Shaded}
\begin{Highlighting}[]
\FunctionTok{cat}\NormalTok{(}\StringTok{"3. IMPLICACIONES PARA LA INTERVENCIÓN INSTITUCIONAL:}\SpecialCharTok{\textbackslash{}n}\StringTok{"}\NormalTok{)}
\end{Highlighting}
\end{Shaded}

\begin{verbatim}
## 3. IMPLICACIONES PARA LA INTERVENCIÓN INSTITUCIONAL:
\end{verbatim}

\begin{Shaded}
\begin{Highlighting}[]
\FunctionTok{cat}\NormalTok{(}\StringTok{"   • Permite identificación de titulaciones en riesgo}\SpecialCharTok{\textbackslash{}n}\StringTok{"}\NormalTok{)}
\end{Highlighting}
\end{Shaded}

\begin{verbatim}
##    • Permite identificación de titulaciones en riesgo
\end{verbatim}

\begin{Shaded}
\begin{Highlighting}[]
\FunctionTok{cat}\NormalTok{(}\StringTok{"   • Facilita diseño de intervenciones personalizadas}\SpecialCharTok{\textbackslash{}n}\StringTok{"}\NormalTok{)}
\end{Highlighting}
\end{Shaded}

\begin{verbatim}
##    • Facilita diseño de intervenciones personalizadas
\end{verbatim}

\begin{Shaded}
\begin{Highlighting}[]
\FunctionTok{cat}\NormalTok{(}\StringTok{"4. CALIDAD DEL MODELO:}\SpecialCharTok{\textbackslash{}n}\StringTok{"}\NormalTok{)}
\end{Highlighting}
\end{Shaded}

\begin{verbatim}
## 4. CALIDAD DEL MODELO:
\end{verbatim}

\begin{Shaded}
\begin{Highlighting}[]
\FunctionTok{cat}\NormalTok{(}\FunctionTok{sprintf}\NormalTok{(}\StringTok{"   • Coeficiente de Silhueta: \%.4f}\SpecialCharTok{\textbackslash{}n}\StringTok{"}\NormalTok{, }\FunctionTok{mean}\NormalTok{(silhouette\_coef[, }\StringTok{"sil\_width"}\NormalTok{])))}
\end{Highlighting}
\end{Shaded}

\begin{verbatim}
##    • Coeficiente de Silhueta: 0.5382
\end{verbatim}

\begin{Shaded}
\begin{Highlighting}[]
\FunctionTok{cat}\NormalTok{(}\StringTok{"   • Separación: Aceptable}\SpecialCharTok{\textbackslash{}n}\StringTok{"}\NormalTok{)}
\end{Highlighting}
\end{Shaded}

\begin{verbatim}
##    • Separación: Aceptable
\end{verbatim}

\begin{Shaded}
\begin{Highlighting}[]
\FunctionTok{cat}\NormalTok{(}\StringTok{"   • Coherencia interna: Buena}\SpecialCharTok{\textbackslash{}n}\StringTok{"}\NormalTok{)}
\end{Highlighting}
\end{Shaded}

\begin{verbatim}
##    • Coherencia interna: Buena
\end{verbatim}

\end{document}
